\documentclass[11pt,letterpaper]{article}
\usepackage{physical_ai_legislation}
\usepackage[backend=biber,style=numeric,sorting=none]{biblatex}
\addbibresource{all_documents.bib}

% Redefine \part to avoid blank pages - just a large centered heading
\makeatletter
\newcounter{docpart}
\renewcommand\part{%
  \@afterindentfalse
  \secdef\@part\@spart}
\renewcommand\@part[2][]{%
  \stepcounter{docpart}%
  \addcontentsline{toc}{part}{#2}%
  \markboth{}{}%
  {\centering
   \interlinepenalty \@M
   \normalfont
   \huge\bfseries #2\par}%
  \nobreak
  \vskip 1em
  \@afterheading}
\renewcommand\@spart[1]{%
  {\centering
   \interlinepenalty \@M
   \normalfont
   \huge\bfseries #1\par}%
  \nobreak
  \vskip 1em
  \@afterheading}
\makeatother

% Fix duplicate hyperref anchors by giving unique names per part
\renewcommand{\theHsection}{\arabic{docpart}.\arabic{section}}
\renewcommand{\theHsubsection}{\arabic{docpart}.\arabic{section}.\arabic{subsection}}

\title{%
  \textbf{Physical AI Oncology Clinical Trial Site}\\
  \textbf{Complete Documentation Package}\\[6pt]
  \large 11 Documents for California's First Physical AI\\
  Oncology Clinical Trial Site\\[6pt]
  \normalsize Legislation, Regulations, Building Code,\\
  Premises Code, and Operations
}
\author{%
  Kevin Kawchak\\
  CEO, ChemicalQDevice\\
  \texttt{kevink@chemicalqdevice.com}
}
\date{March 24, 2026}

\begin{document}
\maketitle
\thispagestyle{fancy}

\begin{center}
\footnotesize
\textit{Unofficial independent draft derived from prior legislation;
no third-party endorsement, sponsorship, affiliation, or authorization
is expressed or implied; and was adapted using Claude Code Opus 4.6.}
\end{center}

\vspace{0.3em}

\begin{center}
\footnotesize
\textit{This work was adapted from original CFR documents in the public domain.
The original ICH document is copyrighted and may be used, reproduced,
incorporated into other works, adapted, modified, translated or distributed
under a public license. This current work is not endorsed or sponsored by
CFR, ICH, or FDA; and was adapted using Claude Code Opus 4.6.}
\end{center}

\vspace{0.3em}

\begin{center}
\footnotesize
\textit{This draft is independent and is not endorsed, sponsored, or approved
by any trial sponsor, CRO, site, IRB, regulator, or medical society;
and was adapted using Claude Code Opus 4.6.}
\end{center}

\vspace{1em}
\tableofcontents


% ============================================================================
% DOCUMENT 1: SB 1042 - California Physical AI Oncology Clinical Trial Authorization and Site Establishment Act of 2026
% ============================================================================

\part{SB 1042 - California Physical AI Oncology Clinical Trial Authorization and Site Establishment Act of 2026}
\setcounter{section}{0}

\section{Legislative Findings and Declarations}

The Legislature finds and declares all of the following:

\begin{enumerate}[label=(\alph*)]
\item Physical artificial intelligence, including robotic surgical systems,
  collaborative robots, radiotherapy positioning systems, imaging assistants,
  rehabilitation exoskeletons, and companion robots, represents a transformative
  advancement in oncology clinical care delivery.

\item A 24-hour on-demand Physical AI oncology clinical trial simulation
  conducted on March 23, 2026, demonstrated that a single facility equipped
  with 10 robot types and 29 robot instances successfully served 168 unique
  patients across 15 cancer types in 24 hours, compared to 5 to 15 patients per
  day in traditional trial formats~\cite{newtrial2026}.

\item The simulation achieved a cumulative site Physical AI Standard Level (PSL)
  score of 63.4 to 64.4, within the Advanced Site band, with 99.7 percent
  facility uptime and an average patient wait time of 8 minutes.

\item All seven adverse events documented during the simulation were Grade 1 or
  Grade 2, all were resolved within the same hour they occurred, and zero
  patient harm resulted, demonstrating that Physical AI systems can manage
  both minute-level clinical details and sustained 24-hour
  operations~\cite{newtrial2026}.

\item The Physical AI trial format shifts control toward the patient by enabling
  patient-chosen scheduling on a 24/7 basis, increasing patient throughput by a
  factor of 10, reducing per-patient costs by 75 to 85 percent, and improving
  retention rates from 70 to 80 percent in traditional formats to 85 to
  95 percent~\cite{formatcomparison2026}.

\item Staffing requirements are reduced from 80 to 120 full-time equivalents in
  traditional trial sites to 8 to 10 full-time equivalents in a Physical AI
  trial site, with the majority of clinical procedures, monitoring,
  documentation, and regulatory compliance tasks performed
  autonomously~\cite{sitespec2026}.

\item The Unification Standard Level (USL) framework provides a standardized
  1.0 to 10.0 scoring system for evaluating physical AI robot readiness across
  four dimensions: simulation framework switching, AI integration, cross-robot
  progress sharing, and clinical trial collaboration~\cite{usl2026}.

\item Three adapted regulatory frameworks, covering ICH E6(R3) good clinical
  practice, 21 CFR Part 50 human subjects protection, and 21 CFR Part 312
  investigational new drug applications, provide comprehensive governance for
  Physical AI oncology trials~\cite{iche6r3adapt,cfr50adapt,cfr312adapt}.

\item The cancer patient journey through a regulated Physical AI oncology trial
  has been documented across 10 sequential stages from pre-screening through
  data lock, demonstrating complete Response outcomes and a reduction in
  recurrence risk from 35 percent to 3 percent~\cite{patientjourney2026}.

\item California has an opportunity to lead the nation in establishing safe,
  regulated, and patient-centric Physical AI oncology clinical trial sites that
  serve more patients with more robots and fewer workers while maintaining the
  highest standards of clinical care.
\end{enumerate}

\section{Definitions}

For purposes of this chapter, the following definitions apply:

\begin{enumerate}[label=(\alph*)]
\item ``Physical AI system'' means a robotic or embodied artificial intelligence
  system that physically interacts with patients, clinical specimens, or medical
  equipment in the context of an oncology clinical trial. Physical AI systems
  include, but are not limited to, the following 10 categories:
  \begin{enumerate}[label=(\arabic*)]
  \item Surgical robots that perform tumor resection through minimally invasive
    procedures.
  \item Collaborative robots (cobots) that stabilize patient limbs and assist
    with needle biopsy procedures.
  \item Radiotherapy patient-positioning robots that align patients with
    submillimeter precision for radiation delivery.
  \item Robotic needle-placement systems that use imaging guidance to place
    biopsy or ablation needles with high accuracy.
  \item Social companion robots that provide emotional support and
    anxiety reduction for patients, including pediatric patients.
  \item Humanoid robots that conduct rehabilitation exercises and patient
    engagement activities.
  \item Radiotherapy motion-management and tracking robots that monitor
    patient breathing and gate radiation beams accordingly.
  \item Imaging assistant robots that perform automated ultrasound and
    diagnostic imaging procedures.
  \item Steerable needle and needle-steering robots that guide flexible
    needles through tissue for targeted ablation.
  \item Rehabilitation exoskeletons and robotic gait trainers that support
    post-surgical mobility recovery.
  \end{enumerate}

\item ``Physical AI Standard Level'' or ``PSL'' means the scoring framework
  that evaluates each robot type on three regulatory dimensions: Omniscient
  (ICH E6(R3) compliance), Omnipresent (21 CFR Part 50 compliance), and
  Omnipotent (21 CFR Part 312 compliance), with each dimension scored from
  0.0 to 10.0 per robot type.

\item ``Unification Standard Level'' or ``USL'' means the 1.0 to 10.0
  scoring framework evaluating physical AI robot readiness across four equally
  weighted dimensions: simulation framework switching, AI integration,
  cross-robot progress sharing, and clinical trial
  collaboration~\cite{usl2026}.

\item ``On-demand trial format'' means a clinical trial structure in which
  patients may schedule procedures at any time during a 24-hour period, 7 days
  per week, with Physical AI systems providing continuous clinical care.

\item ``Digital twin'' means a patient-specific computational model calibrated
  to individual anatomy and treatment parameters, used for procedure planning,
  rehearsal, and intraoperative guidance.

\item ``MCP-PAI standard'' means the Model Context Protocol for Physical AI,
  a five-server interoperability topology providing authorization, FHIR, DICOM,
  ledger, and provenance services with hash-chained audit trails.

\item ``Qualified site'' means a facility that has met all building, premises,
  parking, staffing, equipment, and regulatory requirements established under
  this chapter and has received authorization from the department to conduct
  Physical AI oncology clinical trials.

\item ``Department'' means the California Department of Public Health.

\item ``Sim-to-real validation'' means the process of verifying that a Physical
  AI system's performance in simulated environments transfers to real clinical
  settings within specified tolerances, including trajectory deviation below
  2 millimeters and force discrepancy below 0.5 Newtons.
\end{enumerate}

\section{Authorization to Establish Physical AI\\Oncology Clinical Trial Sites}

\subsection{Site Authorization}

\begin{enumerate}[label=(\alph*)]
\item The department shall establish a process for authorizing Physical AI
  oncology clinical trial sites in California, beginning with the first site
  in a prominent and safe San Francisco location.

\item A qualified entity may apply to the department for authorization to
  establish and operate a Physical AI oncology clinical trial site in a new
  building with dedicated parking facilities.

\item The authorization process shall include review of the following:
  \begin{enumerate}[label=(\arabic*)]
  \item Facility design and construction plans meeting the Physical AI
    oncology trial facility building code standards established under
    this chapter.
  \item Equipment specifications for all 10 Physical AI system categories
    to be deployed at the site.
  \item Staffing plans demonstrating adequate human oversight with a minimum
    of 8 to 10 full-time equivalent positions across three shifts.
  \item USL scores for each Physical AI system, with minimum thresholds of
    7.0 for surgical systems, 6.0 for therapeutic positioning systems, 6.0 for
    diagnostic systems, 4.0 for rehabilitative systems, and 3.0 for companion
    systems.
  \item Phase 0 simulation validation results demonstrating cross-framework
    consistency using two or more simulation frameworks.
  \item Compliance plans for ICH E6(R3), 21 CFR Part 50, and 21 CFR Part 312
    as adapted for Physical AI systems.
  \item Emergency response and cybersecurity plans.
  \item Parking and patient transportation accessibility plans.
  \end{enumerate}

\item The department shall act on a complete application within 120 days of
  receipt.
\end{enumerate}

\subsection{Statewide Applicability}

\begin{enumerate}[label=(\alph*)]
\item The authorization requirements and standards established under this
  chapter shall apply uniformly to Physical AI oncology clinical trial sites
  throughout the State of California.

\item A city, county, or city and county may adopt additional requirements for
  Physical AI oncology clinical trial sites, provided those requirements do not
  conflict with or impose lesser standards than those established under this
  chapter.

\item The department shall publish model guidelines to assist local
  jurisdictions in integrating Physical AI oncology clinical trial site
  requirements into existing land use, building, and health facility
  permitting processes.
\end{enumerate}

\section{Physical AI System Requirements}

\subsection{USL and PSL Compliance}

\begin{enumerate}[label=(\alph*)]
\item Each Physical AI system deployed at an authorized site shall maintain a
  USL score at or above the minimum threshold for its system category as
  specified in Section 3.1(c)(4).

\item The site shall maintain a cumulative PSL score that places it within the
  Advanced Site band (60.0 or above on the 0 to 100 scale) as a condition of
  continued authorization.

\item PSL scores shall be calculated and reported monthly using the three
  regulatory dimensions established in the PSL framework, with changes limited
  to a maximum of 0.3 per hour per dimension per robot type to ensure scoring
  stability~\cite{pslframework2026}.

\item A site whose cumulative PSL score falls below the Advanced Site band for
  two consecutive monthly reporting periods shall submit a corrective action
  plan to the department within 30 days.
\end{enumerate}

\subsection{Simulation Validation Requirements}

\begin{enumerate}[label=(\alph*)]
\item Before deployment at a qualified site, each Physical AI system shall
  complete Phase 0 simulation validation using two or more validated
  frameworks, which may include NVIDIA Isaac Lab, MuJoCo, Gazebo Sim, and
  PyBullet~\cite{iche6r3adapt}.

\item Cross-framework validation shall demonstrate trajectory deviation below
  2 millimeters and force discrepancy below 0.5 Newtons.

\item Digital twin models shall achieve a minimum fidelity of 90 percent when
  compared to patient-specific imaging and anatomical data.

\item Failure mode analysis shall be conducted across a minimum of 1,000
  simulated procedures per Physical AI system type before clinical
  deployment~\cite{cfr312adapt}.
\end{enumerate}

\subsection{Informed Consent and Patient Rights}

\begin{enumerate}[label=(\alph*)]
\item Physical AI oncology clinical trial sites shall comply with all informed
  consent requirements established in the adapted 21 CFR Part 50 for Physical
  AI systems, including disclosure of robotic system roles, USL scores in plain
  language, safety records, and emergency stop provisions~\cite{cfr50adapt}.

\item Patients shall have the right to schedule procedures at any time during
  the 24-hour operating period, consistent with the on-demand trial format.

\item Patients shall have the right to request human intervention or oversight
  at any point during a Physical AI-assisted procedure.

\item Pediatric patients shall receive age-appropriate robotic demonstrations
  and assent procedures as specified in the adapted 21 CFR Part 50
  \S~50.50 through \S~50.56~\cite{cfr50adapt}.

\item Each patient shall be informed of the specific robots that will be
  involved in their care, the role of each robot (autonomous, assistive, or
  shared-autonomy), and the qualifications of human oversight
  personnel~\cite{patientinstructions2026}.
\end{enumerate}

\section{Safety and Reporting Requirements}

\subsection{Adverse Event Reporting}

\begin{enumerate}[label=(\alph*)]
\item The site shall report all Physical AI-related adverse events to the
  department and to the Food and Drug Administration in accordance with the
  timelines established in the adapted 21 CFR Part 312 \S~312.32, including
  7-day reporting for fatal or life-threatening events and 15-day reporting for
  serious events~\cite{cfr312adapt}.

\item Seven categories of reportable Physical AI adverse events shall be
  tracked: serious Physical AI adverse events, Physical AI-drug interaction
  events, cybersecurity incidents, AI model performance degradation,
  sim-to-real divergence, digital twin discrepancy, and safety review committee
  findings.

\item The site shall maintain a Physical AI Safety Review Committee with
  expertise in robotic engineering, AI and machine learning, oncology, and
  cybersecurity, meeting at intervals not exceeding 90 days.
\end{enumerate}

\subsection{Data Protection and Audit Trails}

\begin{enumerate}[label=(\alph*)]
\item All Physical AI systems shall maintain hash-chained audit trails using
  SHA-256 or equivalent cryptographic standards, as required by the adapted
  ICH E6(R3) and 21 CFR Part 11~\cite{iche6r3adapt}.

\item Patient data processed through MCP-PAI servers shall be de-identified
  using HIPAA Safe Harbor methodology with deny-by-default authorization
  policies.

\item The site shall implement federated learning protocols with differential
  privacy protections for multi-site data sharing, as validated in the
  simulation where 70 federated learning rounds across 12 sites used secure
  multi-party computation with differential privacy at epsilon equals
  1.0~\cite{patientjourney2026}.
\end{enumerate}

\section{Facility and Infrastructure Standards}

\begin{enumerate}[label=(\alph*)]
\item Physical AI oncology clinical trial sites shall be constructed in
  accordance with the building code standards established under this chapter,
  including a minimum floor area of 85,000 to 100,000 square feet, HEPA
  filtration in surgical suites, shielded radiotherapy vaults, and dedicated
  server rooms for AI computation infrastructure~\cite{sitespec2026}.

\item The site shall provide dedicated parking facilities meeting the standards
  established under this chapter, including accessible parking, patient
  drop-off zones, and emergency vehicle access.

\item Power infrastructure shall support 800 to 1,200 kilowatts of sustained
  electrical load with redundant power supply and uninterruptible power systems
  for all critical clinical areas.

\item Network infrastructure shall provide a minimum of 10 gigabits per second
  internal bandwidth with cybersecurity protections meeting the requirements of
  the MCP-PAI standard~\cite{sitespec2026}.
\end{enumerate}

\section{Workforce Transition and Training}

\begin{enumerate}[label=(\alph*)]
\item The department, in consultation with the Employment Development Department
  and the California Workforce Development Board, shall establish programs to
  assist clinical trial workers in transitioning to roles in Physical AI trial
  sites or related fields.

\item Each authorized site shall provide comprehensive training for all human
  staff in Physical AI system operation, emergency response, patient
  interaction, and regulatory compliance.

\item Training requirements shall include demonstrated competency in
  understanding Physical AI system capabilities, limitations, and failure modes,
  as well as the ability to intervene manually when necessary.

\item The Legislature recognizes that while Physical AI trial sites require
  fewer workers (8 to 10 full-time equivalents compared to 80 to 120 in
  traditional formats), the transition enables more patients to be treated,
  with the simulation demonstrating 168 patients served in 24 hours
  compared to approximately 5 to 15 per day traditionally.
\end{enumerate}

\section{Relationship to Existing California Law}

This chapter is intended to complement and operate in conjunction with
existing California laws governing health care and artificial intelligence,
including but not limited to the following:

\begin{enumerate}[label=(\alph*)]
\item AB 489 (2025), concerning health care professions and deceptive terms
  related to artificial intelligence, which ensures clear communication about
  AI involvement in clinical care.

\item AB 3030 (2024), concerning health care services and artificial
  intelligence, which establishes transparency requirements for AI in health
  care delivery.

\item SB 1120 (2024), concerning health care coverage and utilization review,
  which addresses AI-assisted clinical decision-making.

\item SB 243 (2025), concerning companion chatbots, which establishes safety
  standards for AI companion systems relevant to social companion robots used
  in Physical AI oncology trials.

\item AB 2013 (2024), concerning generative artificial intelligence and
  training data transparency, which applies to the generative AI and VLA
  models integrated into Physical AI systems.
\end{enumerate}

\section{Appropriation}

The sum of fifty million dollars (\$50,000,000) is hereby appropriated from the
General Fund to the department for the purposes of implementing this chapter,
including the establishment of the site authorization process, development of
model guidelines, workforce transition programs, and ongoing oversight of
authorized Physical AI oncology clinical trial sites.

\section{Effective Date}

This act shall take effect on January 1, 2027.


% ============================================================================
% DOCUMENT 2: AB 2847 - California Physical AI Patient Rights and Robotic Safety Act of 2026
% ============================================================================

\part{AB 2847 - California Physical AI Patient Rights and Robotic Safety Act of 2026}
\setcounter{section}{0}

\section{Legislative Findings and Declarations}

The Legislature finds and declares all of the following:

\begin{enumerate}[label=(\alph*)]
\item Physical AI oncology clinical trials represent a fundamental shift in the
  relationship between patients and clinical care systems, in which control has
  now shifted towards the patient's side, with more patient rights than
  traditional clinical trial formats provide.

\item The 24-hour on-demand Physical AI oncology trial simulation demonstrated
  that patients can choose their own scheduling across a full 24-hour period,
  experience average wait times of 8 minutes, and achieve patient satisfaction
  scores of 4.7 out of 5.0, compared to the rigid scheduling and extended
  wait times typical of traditional oncology trials~\cite{newtrial2026}.

\item More patients will be treated with more robots and fewer workers, with the
  simulation demonstrating 168 patients served by 29 robot instances staffed
  by 8 to 10 human full-time equivalents across three shifts, compared to 5 to
  15 patients per day in traditional formats requiring 80 to 120
  staff~\cite{formatcomparison2026}.

\item The cancer patient journey through a regulated Physical AI oncology trial
  documented a complete patient trajectory from pre-screening through data
  lock across 10 stages, achieving complete Response, R0 resection, and
  reduction in recurrence risk from 35 percent to 3 percent, demonstrating
  that patient outcomes improve under Physical AI
  care~\cite{patientjourney2026}.

\item Pediatric patients require special protections, including age-appropriate
  demonstrations of robotic systems, parental consent provisions, and dedicated
  companion robot interactions, as established in the adapted 21 CFR Part 50
  \S~50.50 through \S~50.56~\cite{cfr50adapt}.

\item Patient-facing robot instructions for all 10 robot types have been
  developed and validated, providing step-by-step guidance for patients
  interacting with surgical robots, cobots, radiotherapy positioning systems,
  needle-placement systems, social companions, humanoids, motion-tracking
  systems, imaging assistants, steerable needles, and rehabilitation
  exoskeletons~\cite{patientinstructions2026}.

\item The simulation demonstrated that both specific details by the minute and
  details over 24 hours were handled correctly by AI, with minute-level vital
  sign monitoring, real-time adverse event detection and resolution, and
  sustained clinical operations across all 24 hours.

\item Seven adverse events occurred across 168 patients (4.2 percent incidence
  rate), all Grade 1 or Grade 2, all resolved within the same hour,
  demonstrating a safety profile that supports expanding patient access to
  Physical AI oncology care~\cite{newtrial2026}.
\end{enumerate}

\section{Patient Bill of Rights for Physical AI\\Oncology Clinical Trials}

\subsection{Right to On-Demand Scheduling}

\begin{enumerate}[label=(\alph*)]
\item Every patient enrolled in a Physical AI oncology clinical trial at an
  authorized California site shall have the right to schedule procedures,
  consultations, and follow-up visits at any time during the site's 24-hour
  operating period, 7 days per week, subject only to clinical safety
  constraints and equipment availability.

\item The site shall provide scheduling systems that enable patients to select
  appointment times without requiring approval from clinical staff, with the
  Physical AI systems managing patient flow and resource allocation
  autonomously.

\item Maximum wait times from arrival to procedure initiation shall not exceed
  30 minutes under normal operating conditions. The simulation demonstrated
  an average wait time of 8 minutes across 168
  patients~\cite{newtrial2026}.
\end{enumerate}

\subsection{Right to Full Disclosure}

\begin{enumerate}[label=(\alph*)]
\item Before any Physical AI-assisted procedure, the patient shall receive clear
  and complete disclosure of the following:
  \begin{enumerate}[label=(\arabic*)]
  \item The specific Physical AI system or systems that will be involved in
    their care, identified by type, model, and the role each system will
    perform (primary, secondary, or ancillary).
  \item Whether the system will operate in autonomous, assistive, or
    shared-autonomy mode.
  \item The system's current USL score presented in plain language, with an
    explanation of what the score means for the patient's
    procedure~\cite{usl2026}.
  \item The system's safety record, including any prior adverse events involving
    the specific system instance.
  \item Emergency stop provisions and how they will be activated if needed.
  \item The identity and qualifications of the human oversight personnel
    responsible for supervising the procedure.
  \item Simulation validation results relevant to the patient's specific
    procedure type.
  \item Digital twin usage and the patient's right to request deletion of their
    digital twin data at any time.
  \end{enumerate}

\item Disclosure materials shall be available in the patient's preferred
  language and in formats accessible to patients with visual, auditory, or
  cognitive impairments.

\item The site shall provide patient-facing robot instructions consistent with
  those developed for each of the 10 Physical AI system
  categories~\cite{patientinstructions2026}.
\end{enumerate}

\subsection{Right to Human Intervention}

\begin{enumerate}[label=(\alph*)]
\item A patient shall have the right to request human intervention or oversight
  at any point during a Physical AI-assisted procedure, and the site shall be
  staffed to accommodate such requests within 2 minutes of the request.

\item A patient may withdraw consent for Physical AI involvement at any time and
  request that the remainder of their procedure be completed by human
  clinicians, to the extent that such transition is clinically safe.

\item The site shall maintain at least one on-call emergency physician and one
  on-call pharmacist at all times during the 24-hour operating period, in
  addition to the site safety officers present on each shift.

\item Emergency stop mechanisms shall be accessible to both clinical staff and
  patients where clinically appropriate, with clear instructions provided in
  advance.
\end{enumerate}

\subsection{Right to Privacy and Data Control}

\begin{enumerate}[label=(\alph*)]
\item All patient data collected, processed, or stored by Physical AI systems
  shall be protected in accordance with HIPAA, the California Consumer Privacy
  Act, the California Privacy Rights Act, and the data governance requirements
  of the adapted ICH E6(R3)~\cite{iche6r3adapt}.

\item Patient data processed through MCP-PAI servers shall use deny-by-default
  authorization policies, meaning no data access is permitted unless
  explicitly authorized for a specific purpose~\cite{iche6r3adapt}.

\item Patients shall have the right to request a complete record of all
  Physical AI system interactions involving their care, including robotic
  telemetry logs, AI decision logs, and digital twin state data.

\item Patients shall have the right to request deletion of their digital twin
  and associated patient-specific data from Physical AI system local storage,
  subject to regulatory data retention requirements under 21 CFR Part 312.

\item Hash-chained audit trails (SHA-256) shall be maintained for all patient
  data transactions, providing an immutable record of data access and
  modification~\cite{iche6r3adapt}.
\end{enumerate}

\subsection{Right to Equitable Access}

\begin{enumerate}[label=(\alph*)]
\item Physical AI oncology clinical trial sites shall not discriminate in
  enrollment, scheduling, or treatment on the basis of race, ethnicity, sex,
  gender identity, sexual orientation, age, disability, socioeconomic status,
  insurance status, or geographic location.

\item The on-demand 24/7 scheduling format shall accommodate patients who work
  non-traditional hours, have caregiving responsibilities, or face
  transportation barriers, consistent with the patient-centric design
  demonstrated in the simulation~\cite{formatcomparison2026}.

\item The site shall provide transportation assistance, including accessible
  parking, patient drop-off zones, and coordination with public transit and
  ride-share services.

\item The per-patient cost reduction of 75 to 85 percent demonstrated in the
  simulation format comparison shall be reflected in reduced financial barriers
  for patients~\cite{formatcomparison2026}.
\end{enumerate}

\subsection{Pediatric Patient Protections}

\begin{enumerate}[label=(\alph*)]
\item Pediatric patients (under 18 years of age) shall receive age-appropriate
  explanations and demonstrations of all Physical AI systems involved in their
  care, consistent with the assent requirements of the adapted 21 CFR Part 50
  \S~50.50 through \S~50.56~\cite{cfr50adapt}.

\item Social companion robots used with pediatric patients shall follow the
  interaction protocols established in the patient robot instructions,
  including gentle high-fives, interactive games, calm-breathing prompts, and
  virtual sticker rewards~\cite{patientinstructions2026}.

\item Humanoid robots conducting rehabilitation exercises with pediatric
  patients shall follow established protocols including slow count-in and
  count-out breathing, posture cues, and interaction durations appropriate for
  the child's age and condition~\cite{patientinstructions2026}.

\item A parent or guardian shall be present or immediately available during all
  Physical AI-assisted procedures involving pediatric patients, and an advocate
  shall be appointed for wards of the state consistent with \S~50.56 of the
  adapted regulations~\cite{cfr50adapt}.

\item The simulation demonstrated successful management of pediatric patients
  including pediatric ALL (acute lymphoblastic leukemia) and AML (acute
  myeloid leukemia) cases, with companion robot-assisted anxiety reduction
  achieving 95 percent effectiveness~\cite{newtrial2026}.
\end{enumerate}

\section{Robotic Safety Standards}

\subsection{Physical AI System Safety Certification}

\begin{enumerate}[label=(\alph*)]
\item Before deployment at an authorized California site, each Physical AI
  system shall obtain safety certification through a process administered by
  the department, demonstrating the following:
  \begin{enumerate}[label=(\arabic*)]
  \item Completion of Phase 0 simulation validation across two or more
    validated simulation frameworks with cross-framework trajectory deviation
    below 2 millimeters and force discrepancy below 0.5
    Newtons~\cite{iche6r3adapt}.
  \item Achievement of minimum USL scores: 7.0 for surgical robots, 6.0 for
    therapeutic positioning and diagnostic systems, 4.0 for rehabilitative
    systems, and 3.0 for companion systems~\cite{cfr50adapt}.
  \item Completion of failure mode analysis across a minimum of 1,000
    simulated procedures~\cite{cfr312adapt}.
  \item Functional emergency stop mechanisms tested and verified.
  \item Force-limiting capabilities appropriate to the system type, including
    maximum 250 millimeters per second approach speed and force-limited contact
    for cobots~\cite{patientinstructions2026}.
  \end{enumerate}

\item Safety certification shall be renewed annually and following any
  significant hardware or software modification.
\end{enumerate}

\subsection{Continuous Safety Monitoring}

\begin{enumerate}[label=(\alph*)]
\item Each Physical AI system shall maintain continuous telemetry recording
  during all patient interactions, as demonstrated in the hourly robot logs of
  the 24-hour simulation where all 29 robot instances were tracked with
  active and standby status, performance metrics, and state
  transitions~\cite{newtrial2026}.

\item Real-time safety monitoring shall include automated detection of force
  limit exceedances, positioning errors beyond specified tolerances, AI model
  performance degradation, and cybersecurity anomalies.

\item The simulation demonstrated 99.7 percent fleet uptime across 29 robot
  instances over 24 hours, with one planned maintenance event (SURG-01
  maintenance window from 22:15 to 01:59) managed without patient
  impact~\cite{newtrial2026}.

\item Sim-to-real divergence exceeding 4 millimeters in trajectory or 1.0
  Newton in force shall trigger an automatic safety review and potential
  suspension of the affected system~\cite{cfr312adapt}.
\end{enumerate}

\subsection{Adverse Event Response}

\begin{enumerate}[label=(\alph*)]
\item Every Physical AI-related adverse event shall be documented, reported, and
  resolved in accordance with the timelines and categories established in the
  adapted 21 CFR Part 312 \S~312.32~\cite{cfr312adapt}.

\item The site shall achieve a target of resolving all Grade 1 and Grade 2
  adverse events within the same hour of occurrence, consistent with the
  performance demonstrated in the simulation where all seven adverse events
  (ranging from nausea and hypotension to pain, cough, desaturation, and
  anxiety) were resolved within the hour~\cite{newtrial2026}.

\item A Physical AI Safety Review Committee shall convene within 24 hours of any
  Grade 3 or higher adverse event and within 7 days of any fatal or
  life-threatening event.

\item The committee shall include expertise in robotic engineering, AI and
  machine learning, oncology, and cybersecurity, meeting at regular intervals
  not exceeding 90 days~\cite{cfr312adapt}.
\end{enumerate}

\section{Patient Complaint and Grievance Process}

\begin{enumerate}[label=(\alph*)]
\item Each authorized site shall establish a patient complaint and grievance
  process that is accessible 24 hours per day, 7 days per week, consistent
  with the site's operating schedule.

\item Complaints involving Physical AI system malfunctions or safety concerns
  shall be investigated within 48 hours and the patient shall be notified of
  the investigation outcome within 14 days.

\item The department shall maintain a public registry of Physical AI-related
  patient complaints and their resolutions, with appropriate de-identification
  of patient information.

\item Patients shall have the right to file complaints directly with the
  department if they believe the site's internal grievance process has not
  adequately addressed their concerns.
\end{enumerate}

\section{Enforcement}

\begin{enumerate}[label=(\alph*)]
\item The department shall have authority to inspect authorized Physical AI
  oncology clinical trial sites without prior notice to verify compliance
  with this article.

\item Violations of patient rights established under this article shall be
  subject to civil penalties of not less than one thousand dollars (\$1,000)
  and not more than twenty-five thousand dollars (\$25,000) per violation.

\item The department may suspend or revoke site authorization for repeated or
  serious violations of patient rights or robotic safety standards.

\item Nothing in this article shall be construed to limit any other remedy
  available to patients under California law.
\end{enumerate}

\section{Effective Date}

This act shall take effect on January 1, 2027.


% ============================================================================
% DOCUMENT 3: SB 892 - California Physical AI Clinical Data Protection and Transparency Act of 2026
% ============================================================================

\part{SB 892 - California Physical AI Clinical Data Protection and Transparency Act of 2026}
\setcounter{section}{0}

\section{Legislative Findings and Declarations}

The Legislature finds and declares all of the following:

\begin{enumerate}[label=(\alph*)]
\item Physical AI oncology clinical trials generate unprecedented volumes of
  patient data, including robotic telemetry, AI decision logs, digital twin
  state data, medical imaging, vital signs captured at minute-level resolution,
  and federated learning model parameters~\cite{newtrial2026}.

\item The 24-hour on-demand simulation generated 178 output files across 25
  sequential commits, including hourly patient records with minute-level vital
  sign data for 168 patients, detailed robot telemetry for 29 instances across
  24 hours, and PSL scoring data for all 10 robot
  types~\cite{newtrial2026}.

\item The MCP-PAI standard establishes a five-server interoperability topology
  (authorization, FHIR, DICOM, ledger, provenance) with 23 tools, six actor
  roles, five conformance levels, and hash-chained audit trails, providing a
  comprehensive data governance architecture~\cite{iche6r3adapt}.

\item Training data transparency requirements established by AB 2013 (2024)
  apply to the generative AI and VLA models integrated into Physical AI
  systems, and additional transparency requirements specific to robotic
  telemetry and clinical decision data are needed.

\item The cancer patient journey demonstrated federated learning across 12 sites
  using secure multi-party computation with differential privacy at epsilon
  equals 1.0, providing a model for privacy-preserving multi-site data
  sharing~\cite{patientjourney2026}.

\item Both specific details by the minute and details over 24 hours were
  captured and managed correctly by AI systems, creating a data governance
  challenge that requires clear rules for data retention, access, sharing, and
  deletion across multiple timescales.
\end{enumerate}

\section{Definitions}

For purposes of this chapter, the following definitions apply:

\begin{enumerate}[label=(\alph*)]
\item ``Physical AI clinical data'' means any data generated, collected,
  processed, or stored by Physical AI systems during oncology clinical trials,
  including robotic telemetry data, AI decision logs, digital twin state data,
  patient vital signs, medical imaging data, procedural records, adverse event
  records, and PSL and USL scoring data.

\item ``Robotic telemetry data'' means real-time operational data from Physical
  AI systems including position, force, velocity, temperature, power
  consumption, active or standby status, and performance metrics.

\item ``AI decision log'' means a record of decisions made by AI algorithms
  integrated into Physical AI systems, including treatment recommendations,
  scheduling decisions, resource allocation, and safety interventions.

\item ``Digital twin data'' means patient-specific computational model data
  including anatomical models, treatment simulation parameters, procedure
  rehearsal records, and predictive model outputs.

\item ``Hash-chained audit trail'' means an immutable, cryptographically linked
  sequence of data access and modification records using SHA-256 or equivalent
  algorithms, as specified in the adapted ICH E6(R3)~\cite{iche6r3adapt}.

\item ``Federated learning data'' means model parameters, gradient updates, and
  aggregated statistics shared between Physical AI systems across multiple
  clinical trial sites using privacy-preserving computation methods.

\item ``De-identification'' means the process of removing or obscuring personal
  health information using the HIPAA Safe Harbor methodology with 18
  identifier categories, as required by the adapted regulatory
  frameworks~\cite{cfr50adapt}.
\end{enumerate}

\section{Data Collection Transparency}

\subsection{Notice Requirements}

\begin{enumerate}[label=(\alph*)]
\item Before enrollment, each patient shall receive a comprehensive data
  collection notice specifying all categories of Physical AI clinical data that
  will be generated during their participation, the purposes for which each
  category will be used, and the retention period for each category.

\item The notice shall identify the five MCP-PAI servers that will process
  patient data (authorization, FHIR, DICOM, ledger, provenance) and describe
  the role of each server in plain language~\cite{iche6r3adapt}.

\item The notice shall disclose whether patient data will be used for federated
  learning and, if so, describe the differential privacy protections applied,
  including the privacy budget (epsilon value)~\cite{patientjourney2026}.

\item The notice shall be updated whenever new categories of data collection are
  introduced or existing categories are materially changed.
\end{enumerate}

\subsection{Real-Time Data Access}

\begin{enumerate}[label=(\alph*)]
\item Patients shall have the right to access their Physical AI clinical data
  in real time through a secure patient portal, including current vital signs,
  procedure status, robot telemetry relevant to their care, and digital twin
  visualizations.

\item The site shall provide data access in standardized formats, including
  FHIR for clinical data and DICOM for medical imaging data.

\item Data access requests shall be fulfilled within 24 hours for historical
  data not available through the real-time portal.
\end{enumerate}

\section{Data Protection Standards}

\subsection{Encryption and Access Control}

\begin{enumerate}[label=(\alph*)]
\item All Physical AI clinical data shall be encrypted at rest using AES-256 or
  equivalent and in transit using TLS 1.3 or later.

\item Access to Physical AI clinical data shall follow deny-by-default
  authorization policies, meaning no access is permitted unless explicitly
  authorized for a specific purpose, role, and time period~\cite{iche6r3adapt}.

\item The six MCP-PAI actor roles (patient, clinician, investigator, sponsor,
  regulator, system) shall define the maximum data access scope for each
  role, with no role having unrestricted access to all data
  categories~\cite{iche6r3adapt}.

\item Multi-factor authentication shall be required for all human access to
  Physical AI clinical data systems.
\end{enumerate}

\subsection{Audit Trail Requirements}

\begin{enumerate}[label=(\alph*)]
\item Hash-chained audit trails shall be maintained for all Physical AI clinical
  data transactions, providing cryptographic proof of data integrity and an
  immutable record of who accessed what data, when, and for what
  purpose~\cite{iche6r3adapt}.

\item Audit trail records shall be retained for a minimum of 15 years following
  the completion of the clinical trial, or longer if required by applicable
  regulations.

\item The site shall conduct quarterly audit trail integrity verification using
  cryptographic chain validation.

\item Audit trail data shall be stored in a separate system from clinical data
  to prevent simultaneous compromise.
\end{enumerate}

\subsection{De-identification Standards}

\begin{enumerate}[label=(\alph*)]
\item All Physical AI clinical data shared outside the treating site shall be
  de-identified using the HIPAA Safe Harbor methodology, removing all 18
  categories of identifiers~\cite{cfr50adapt}.

\item Robotic telemetry data shall be assessed for re-identification risk
  before sharing, as movement patterns and procedural signatures may
  constitute quasi-identifiers.

\item Digital twin data shall be fully de-identified before any use in
  federated learning, research, or quality improvement activities.

\item The site shall employ a qualified statistician or privacy officer to
  verify that de-identified data sets meet the threshold for non-identifiability.
\end{enumerate}

\section{Federated Learning and Multi-Site\\Data Sharing}

\begin{enumerate}[label=(\alph*)]
\item Physical AI oncology clinical trial sites that participate in federated
  learning shall implement the privacy-preserving computation methods
  demonstrated in the cancer patient journey, including secure multi-party
  computation with differential privacy~\cite{patientjourney2026}.

\item The privacy budget (epsilon value) for differential privacy shall not
  exceed 1.0 per federated learning round without explicit patient consent and
  IRB approval.

\item Model parameters shared through federated learning shall be verified to
  not contain memorized patient data through membership inference testing.

\item Federated learning participants shall maintain data harmonization
  standards using DICOM, FHIR, ICD-10 to SNOMED CT mapping, and LOINC
  coding~\cite{iche6r3adapt}.

\item The simulation demonstrated that 70 federated learning rounds across 12
  sites achieved clinically significant model improvements while maintaining
  patient privacy, providing evidence that multi-site Physical AI data sharing
  can be conducted safely~\cite{patientjourney2026}.
\end{enumerate}

\section{Data Retention and Deletion}

\begin{enumerate}[label=(\alph*)]
\item Physical AI clinical data shall be retained for the duration required by
  applicable federal and state regulations, including a minimum of 2 years
  following investigational drug approval or 2 years after investigation
  discontinuation, consistent with 21 CFR Part 312~\cite{cfr312adapt}.

\item Upon trial completion and satisfaction of regulatory retention
  requirements, patients shall have the right to request deletion of their
  digital twin data and patient-specific data from Physical AI system local
  storage.

\item Deletion requests shall be fulfilled within 30 days using secure data
  erasure methods that comply with HIPAA Safe Harbor requirements, and the
  site shall provide written confirmation of deletion~\cite{cfr312adapt}.

\item Aggregated, de-identified data may be retained indefinitely for research
  purposes, provided it meets the de-identification standards established in
  Section 4.3.

\item When a site's authorization is revoked or a trial is terminated, all
  patient-specific data shall be securely erased from Physical AI system local
  storage within 30 days of decommissioning, consistent with the IND
  withdrawal requirements of the adapted 21 CFR Part 312
  \S~312.38~\cite{cfr312adapt}.
\end{enumerate}

\section{Cybersecurity Incident Reporting}

\begin{enumerate}[label=(\alph*)]
\item Any cybersecurity incident affecting Physical AI clinical data shall be
  reported to the department, to affected patients, and to the FDA in
  accordance with the following timelines:
  \begin{enumerate}[label=(\arabic*)]
  \item Within 7 days if the incident involves potential patient harm or
    breach of protected health information for 500 or more individuals.
  \item Within 15 days for all other cybersecurity incidents affecting Physical
    AI systems~\cite{cfr312adapt}.
  \end{enumerate}

\item The site shall maintain a cybersecurity incident response plan that
  addresses Physical AI-specific threats including robot control system
  compromise, AI model poisoning, digital twin data manipulation, and
  MCP-PAI server infrastructure attacks.

\item Annual cybersecurity assessments shall include penetration testing of
  Physical AI system interfaces, MCP-PAI server infrastructure, and network
  segmentation between clinical and administrative systems.

\item The simulation demonstrated continuous operation across 24 hours with no
  cybersecurity incidents, including secure management of 10 gigabits per
  second internal network traffic supporting 29 robot instances and patient
  data systems~\cite{sitespec2026}.
\end{enumerate}

\section{AI Model Transparency}

\begin{enumerate}[label=(\alph*)]
\item In accordance with the principles of AB 2013 (2024), sites shall maintain
  and make available upon request documentation of the AI models integrated
  into Physical AI systems, including model architecture, training data
  sources, validation methodology, and performance benchmarks.

\item The five categories of AI recognized in the adapted ICH E6(R3) -
  generative AI (VLA models, diffusion policies), agentic AI (LLM-based
  assistants), reinforcement learning, self-supervised learning, and
  supervised learning - shall each have category-specific transparency
  requirements~\cite{iche6r3adapt}.

\item Pre-determined change control plans (PCCPs) for AI model updates shall be
  filed with the department and made available to patients upon request,
  specifying which model changes require protocol amendments versus those
  within the PCCP scope~\cite{cfr312adapt}.

\item AI model performance shall be reported in the site's annual report
  alongside PSL and USL scores, providing a comprehensive transparency
  record~\cite{cfr312adapt}.
\end{enumerate}

\section{Enforcement and Penalties}

\begin{enumerate}[label=(\alph*)]
\item Violations of data protection standards established under this chapter
  shall be subject to civil penalties of not less than two thousand five
  hundred dollars (\$2,500) and not more than seven thousand five hundred
  dollars (\$7,500) per violation per patient affected.

\item Intentional or reckless failure to maintain hash-chained audit trails or
  to report cybersecurity incidents within required timelines shall constitute
  an aggravated violation subject to penalties of up to twenty-five thousand
  dollars (\$25,000) per violation.

\item The Attorney General may bring an action to enforce this chapter on behalf
  of the people of the state.

\item Nothing in this chapter shall be construed to limit any private right of
  action available to patients under the California Consumer Privacy Act or
  the California Privacy Rights Act.
\end{enumerate}

\section{Effective Date}

This act shall take effect on January 1, 2027.


% ============================================================================
% DOCUMENT 4: San Francisco Municipal Code Update - Physical AI Oncology Clinical Trial Site Requirements
% ============================================================================

\part{San Francisco Municipal Code Update - Physical AI Oncology Clinical Trial Site Requirements}
\setcounter{section}{0}

\section{Purpose and Scope}

\subsection{Purpose}

This regulation establishes the municipal permit, zoning, and operational
requirements for Physical AI oncology clinical trial sites within the City and
County of San Francisco, in furtherance of the California Physical AI Oncology
Clinical Trial Authorization and Site Establishment Act (SB 1042). San
Francisco shall be the location of the first Physical AI oncology clinical trial
site in California, situated in a prominent and safe location that provides
accessible care to oncology patients throughout the city and surrounding Bay
Area communities.

\subsection{Scope}

These requirements apply to the permitting, construction, and operation of
new-construction facilities specifically designed as Physical AI oncology
clinical trial sites, including the associated parking facilities required under
this regulation. These requirements supplement but do not replace existing San
Francisco Planning Code, Building Code, and Health Code requirements for medical
facilities.

\subsection{Evidence Base}

The 24-hour on-demand Physical AI oncology clinical trial simulation conducted
on March 23, 2026, provides the primary evidence base for these regulations,
demonstrating that a single 85,000 to 100,000 square-foot facility equipped
with 10 robot types (29 total instances) served 168 unique patients across 15
cancer types with 99.7 percent facility uptime and zero patient
harm~\cite{newtrial2026}. The simulation demonstrated that both specific
clinical details at minute-level resolution and sustained operations across a
full 24-hour cycle were managed correctly by Physical AI
systems~\cite{newtrial2026}.

\section{Zoning Requirements}

\subsection{Permitted Zones}

\begin{enumerate}[label=(\alph*)]
\item Physical AI oncology clinical trial sites shall be a conditionally
  permitted use in the following San Francisco zoning districts:
  \begin{enumerate}[label=(\arabic*)]
  \item Neighborhood Commercial (NC) districts with parcels of 20,000 square
    feet or larger.
  \item Mixed Use-General (MUG) districts.
  \item Production, Distribution, and Repair (PDR) districts adjacent to
    residential or commercial zones with transit access.
  \item Downtown General (DTG) and Downtown Office (C-3) districts.
  \item Medical Use Districts and Health Care Services Districts.
  \end{enumerate}

\item Preference shall be given to locations within one-quarter mile of a Bay
  Area Rapid Transit (BART) station or San Francisco Municipal Railway (Muni)
  stop to maximize patient accessibility consistent with the on-demand 24/7
  scheduling model~\cite{formatcomparison2026}.

\item The site shall be located in a prominent and safe location, defined as a
  parcel with frontage on a street classified as an arterial or collector in
  the San Francisco General Plan, with adequate street lighting, pedestrian
  infrastructure, and proximity to emergency services.
\end{enumerate}

\subsection{Setback and Height Requirements}

\begin{enumerate}[label=(\alph*)]
\item The minimum lot size for a Physical AI oncology clinical trial site shall
  be 40,000 square feet, accommodating the 85,000 to 100,000 square-foot
  building footprint across multiple stories plus dedicated parking
  facilities~\cite{sitespec2026}.

\item Building height shall comply with the applicable district height limit
  but shall not exceed 65 feet unless a variance is obtained through the
  Planning Commission conditional use process.

\item Street-level setbacks shall provide a minimum 15-foot pedestrian zone
  along all public street frontages.

\item The site shall include a dedicated patient drop-off and pick-up zone
  accessible from the public right-of-way, with a covered canopy extending a
  minimum of 20 feet from the building entrance.
\end{enumerate}

\section{Conditional Use Permit Requirements}

\subsection{Application Contents}

An application for a conditional use permit for a Physical AI oncology clinical
trial site shall include, in addition to standard Planning Code requirements,
the following:

\begin{enumerate}[label=(\alph*)]
\item A Physical AI System Deployment Plan identifying all 10 robot types to be
  deployed, the number of instances of each type, and the USL and PSL scores
  for each system, demonstrating compliance with the minimum thresholds
  established by state law~\cite{usl2026,pslframework2026}.

\item A Patient Flow Analysis based on the simulation data, including projected
  daily patient volumes (up to 168 patients per 24 hours), arrival and
  departure patterns, peak utilization periods (the simulation showed a peak
  of 15 arrivals in hour 09 with 72.4 percent utilization and 28 concurrent
  patients on-site), and measures to manage queuing and
  wait times~\cite{newtrial2026}.

\item A Neighborhood Impact Assessment addressing noise, traffic, parking
  demand, emergency vehicle access, and the visual impact of the facility on
  the surrounding neighborhood.

\item A Community Benefits Plan describing workforce development programs,
  local hiring commitments, and patient access programs for San Francisco
  residents.

\item A Transportation Demand Management (TDM) Plan meeting the requirements
  of San Francisco Planning Code Section 169, reflecting the 24/7 operating
  schedule and the reduced staffing model (8 to 10 full-time equivalents
  versus 80 to 120 in traditional trial sites)~\cite{sitespec2026}.
\end{enumerate}

\subsection{Public Notice and Hearing}

\begin{enumerate}[label=(\alph*)]
\item The Planning Department shall provide notice of the conditional use
  permit application to all owners and occupants of property within 300 feet
  of the proposed site.

\item A public hearing before the Planning Commission shall be held not fewer
  than 30 days after notice is provided.

\item The applicant shall hold at least one community meeting in the
  neighborhood where the site is proposed, at least 14 days before the
  Planning Commission hearing.
\end{enumerate}

\section{Building Permit Requirements}

\subsection{Structural and Mechanical Systems}

\begin{enumerate}[label=(\alph*)]
\item Building permit applications for Physical AI oncology clinical trial
  sites shall demonstrate compliance with the California Building Standards
  Code (Title 24) as supplemented by the San Francisco Building Code, with
  the following additional requirements specific to Physical AI operations:

\item Surgical suites shall incorporate HEPA filtration (ISO Class 7 or better),
  positive-pressure air handling, and floor load capacity of at least 150
  pounds per square foot to support surgical robot systems and associated
  equipment~\cite{sitespec2026}.

\item Radiotherapy vaults shall include radiation shielding meeting the
  requirements of the California Radiologic Health Branch, with additional
  structural reinforcement to support the weight of radiotherapy positioning
  robot systems~\cite{sitespec2026}.

\item Server rooms for AI computation infrastructure shall be located on a
  separate HVAC zone with redundant cooling capable of rejecting a minimum
  heat load of 200 kilowatts.

\item The building shall accommodate 800 to 1,200 kilowatts of sustained
  electrical load with redundant utility feeds, automatic transfer switches,
  and uninterruptible power systems covering all critical clinical and
  data systems~\cite{sitespec2026}.
\end{enumerate}

\subsection{Accessibility}

\begin{enumerate}[label=(\alph*)]
\item The site shall exceed the minimum accessibility requirements of the
  Americans with Disabilities Act and California Building Code Chapter 11B
  to ensure that the on-demand patient scheduling model is fully accessible
  to patients with mobility, visual, auditory, or cognitive impairments.

\item All patient circulation routes shall be designed to accommodate hospital
  beds, wheelchairs, and rehabilitation exoskeletons simultaneously.

\item Wayfinding systems shall include visual, auditory, and tactile elements
  appropriate for a facility where patients interact with 10 categories of
  robotic systems~\cite{patientinstructions2026}.
\end{enumerate}

\section{Health Code Requirements}

\subsection{Operating License}

\begin{enumerate}[label=(\alph*)]
\item A Physical AI oncology clinical trial site shall obtain a San Francisco
  Health Care Facility Operating License in addition to the state authorization
  required under SB 1042.

\item The operating license shall specify the 10 categories of Physical AI
  systems authorized for use at the site, the maximum number of concurrent
  patients, and the operating schedule (24 hours per day, 7 days per week).

\item The license shall be renewed annually, contingent upon maintaining a
  cumulative PSL score within the Advanced Site band (60.0 or above) and
  submitting the annual safety and performance
  reports~\cite{pslframework2026}.
\end{enumerate}

\subsection{Inspection Requirements}

\begin{enumerate}[label=(\alph*)]
\item The San Francisco Department of Public Health shall conduct an initial
  pre-operational inspection of the facility before the first patient is
  enrolled.

\item Routine inspections shall occur at least quarterly, with additional
  inspections triggered by adverse event reports, cybersecurity incidents, or
  patient complaints.

\item Inspections shall include verification of Physical AI system safety
  certifications, review of adverse event logs and resolution records, audit
  trail integrity checks, and observation of at least one active procedure.

\item Inspection findings shall be posted on the department's public website
  within 30 days of the inspection date.
\end{enumerate}

\section{Noise and Environmental Standards}

\begin{enumerate}[label=(\alph*)]
\item The site shall comply with the San Francisco Noise Ordinance (Article 29
  of the Police Code) with the following additional provisions for 24-hour
  Physical AI operations:
  \begin{enumerate}[label=(\arabic*)]
  \item HVAC and mechanical equipment noise at the property line shall not
    exceed 45 dBA during nighttime hours (10:00 PM to 7:00 AM).
  \item Loading and delivery activities shall be restricted to daytime hours
    (7:00 AM to 8:00 PM) except for emergency medical supply deliveries.
  \item Robot maintenance and testing activities generating noise above ambient
    levels shall be conducted in acoustically isolated maintenance bays.
  \end{enumerate}

\item Emergency generator testing shall be limited to daytime hours and shall
  not exceed one test per month lasting no more than 30 minutes.

\item The site shall implement a stormwater management plan compliant with San
  Francisco's Green Infrastructure Ordinance, incorporating the parking
  facility's impervious surface area.
\end{enumerate}

\section{Relationship to State Requirements}

\begin{enumerate}[label=(\alph*)]
\item These municipal requirements supplement and do not replace the
  state-level requirements established under SB 1042, AB 2847, and SB 892.

\item Where a municipal requirement exceeds the corresponding state requirement,
  the municipal requirement shall apply. Where a state requirement exceeds the
  corresponding municipal requirement, the state requirement shall apply.

\item The three adapted regulatory frameworks (ICH E6(R3), 21 CFR Part 50, and
  21 CFR Part 312) apply to Physical AI oncology clinical trial sites in San
  Francisco as they apply throughout the state~\cite{iche6r3adapt,cfr50adapt,cfr312adapt}.
\end{enumerate}


% ============================================================================
% DOCUMENT 5: California Code of Regulations, Title 22 - Physical AI Oncology Trial Site Regulations
% ============================================================================

\part{California Code of Regulations, Title 22 - Physical AI Oncology Trial Site Regulations}
\setcounter{section}{0}

\section{General Provisions}

\subsection{\S~75000. Authority and Purpose}

\begin{enumerate}[label=(\alph*)]
\item These regulations are adopted pursuant to the authority granted to the
  California Department of Public Health (hereinafter ``the department'') under
  the California Physical AI Oncology Clinical Trial Authorization and Site
  Establishment Act of 2026 (Health and Safety Code \S~1399.900 et seq.).

\item The purpose of these regulations is to establish standards and procedures
  for the authorization, inspection, and oversight of Physical AI oncology
  clinical trial sites throughout California.

\item These regulations apply uniformly across the State of California. The
  first authorized site shall be established in San Francisco, and subsequent
  sites may be authorized in any California jurisdiction that meets the
  requirements of these regulations.
\end{enumerate}

\subsection{\S~75001. Definitions}

In addition to the definitions established in Health and Safety Code
\S~1399.901, the following definitions apply:

\begin{enumerate}[label=(\alph*)]
\item ``Authorization certificate'' means the document issued by the department
  to a qualified entity permitting the operation of a Physical AI oncology
  clinical trial site.

\item ``Graduated autonomy'' means the phased approach to Physical AI system
  deployment in which autonomy levels increase from teleoperation in Phase 1
  to supervised autonomy in Phase 2 to full autonomy range in Phase 3, as
  established in the adapted 21 CFR Part 312 \S~312.21~\cite{cfr312adapt}.

\item ``Pre-procedure safety matrix'' means the checklist of safety
  verifications that shall be completed before initiating any Physical
  AI-assisted clinical procedure, as defined in the adapted 21 CFR Part 50
  \S~50.3~\cite{cfr50adapt}.

\item ``Robot capability profile'' means the formal documentation of a Physical
  AI system's validated task capabilities, performance envelope, and safety
  constraints.

\item ``Site safety officer'' means a human staff member certified by the
  department to supervise Physical AI system operations during a shift, with
  authority to initiate emergency stop procedures and override autonomous
  operations.
\end{enumerate}

\section{Authorization Process}

\subsection{\S~75010. Application Requirements}

\begin{enumerate}[label=(\alph*)]
\item An entity seeking authorization to operate a Physical AI oncology
  clinical trial site shall submit an application to the department on forms
  prescribed by the department, including the following:
  \begin{enumerate}[label=(\arabic*)]
  \item Identification of the applicant entity, including ownership structure,
    officers, and any parent organizations.
  \item Site location and facility description, including architectural plans,
    engineering specifications, and evidence of compliance with applicable
    building codes.
  \item Physical AI System Deployment Plan identifying all robot types and
    instances, with current USL scores for each system~\cite{usl2026}.
  \item Phase 0 simulation validation results for each Physical AI system
    type, demonstrating cross-framework validation using two or more
    frameworks~\cite{iche6r3adapt}.
  \item Staffing plan identifying all positions, qualifications, training
    requirements, and shift schedules for the 24-hour operating model.
  \item Quality management system documentation addressing Physical AI-specific
    processes.
  \item Emergency response plan addressing Physical AI system failures,
    cybersecurity incidents, natural disasters, and medical emergencies.
  \item Patient flow projections based on the simulation data showing capacity
    for 150 to 180 patients per day across 15 or more cancer
    types~\cite{newtrial2026}.
  \item Financial documentation demonstrating adequate resources to sustain
    operations for a minimum of 24 months.
  \item Evidence of compliance with adapted ICH E6(R3), 21 CFR Part 50, and
    21 CFR Part 312 regulatory frameworks~\cite{iche6r3adapt,cfr50adapt,cfr312adapt}.
  \end{enumerate}

\item The application shall be accompanied by a nonrefundable filing fee of
  twenty-five thousand dollars (\$25,000).
\end{enumerate}

\subsection{\S~75011. Review Process}

\begin{enumerate}[label=(\alph*)]
\item The department shall assign a review team consisting of specialists in
  robotic engineering, oncology, clinical trial regulation, facility safety,
  and cybersecurity.

\item The review team shall complete an initial completeness review within 30
  days of application receipt and notify the applicant of any deficiencies.

\item The department shall conduct a substantive review of the complete
  application and issue a decision within 120 days of receiving a complete
  application.

\item The department may conduct a pre-authorization site inspection as part
  of the review process.
\end{enumerate}

\subsection{\S~75012. Authorization Certificate}

\begin{enumerate}[label=(\alph*)]
\item An authorization certificate shall specify the following:
  \begin{enumerate}[label=(\arabic*)]
  \item The authorized Physical AI system types and maximum number of instances.
  \item The maximum number of concurrent patients.
  \item The authorized operating schedule.
  \item The minimum staffing requirements.
  \item Any conditions or limitations specific to the site.
  \end{enumerate}

\item The authorization certificate shall be valid for a period of three years,
  subject to annual compliance verification.

\item The certificate shall be prominently displayed at the site entrance and
  available for public inspection.
\end{enumerate}

\section{Physical AI System Standards}

\subsection{\S~75020. System Registration}

\begin{enumerate}[label=(\alph*)]
\item Each Physical AI system instance deployed at an authorized site shall be
  registered with the department, including the following information:
  \begin{enumerate}[label=(\arabic*)]
  \item System type, manufacturer, model, and serial number.
  \item Current USL score with dimension-level breakdown (simulation framework
    switching, AI integration, cross-robot progress sharing, clinical trial
    collaboration)~\cite{usl2026}.
  \item Current PSL scores across three regulatory dimensions (Omniscient,
    Omnipresent, Omnipotent)~\cite{pslframework2026}.
  \item Robot capability profile documenting validated task capabilities.
  \item Phase 0 simulation validation results.
  \item Installation date and commissioning records.
  \end{enumerate}

\item Registration shall be updated within 14 days of any change to system
  configuration, USL score, or capability profile.
\end{enumerate}

\subsection{\S~75021. Minimum USL Thresholds}

Physical AI systems shall meet the following minimum USL score thresholds,
consistent with the MCP-PAI standard~\cite{cfr50adapt}:

\begin{center}
\begin{tabular}{lc}
\toprule
\textbf{System Category} & \textbf{Minimum USL Score} \\
\midrule
Surgical robots & 7.0 \\
Radiotherapy positioning robots & 6.0 \\
Robotic needle-placement systems & 6.0 \\
Radiotherapy motion-tracking robots & 6.0 \\
Imaging assistant robots & 6.0 \\
Steerable needle robots & 6.0 \\
Collaborative robots (cobots) & 5.0 \\
Rehabilitation exoskeletons & 4.0 \\
Humanoid robots & 4.0 \\
Social companion robots & 3.0 \\
\bottomrule
\end{tabular}
\end{center}

\subsection{\S~75022. Maintenance and Calibration}

\begin{enumerate}[label=(\alph*)]
\item Each Physical AI system shall undergo scheduled maintenance at intervals
  specified by the manufacturer and documented in the robot capability profile.

\item Maintenance windows shall be coordinated to maintain continuous clinical
  capability across all system categories, as demonstrated in the simulation
  where SURG-01 maintenance (22:15 to 01:59) was scheduled during low-demand
  hours with SURG-02 and SURG-03 remaining
  available~\cite{newtrial2026}.

\item Calibration records shall include date, calibration parameters, results,
  and the identity of the person or system performing the calibration.

\item A Physical AI system shall not be used for patient procedures until
  post-maintenance verification demonstrates performance within specified
  tolerances.
\end{enumerate}

\section{Staffing and Training Standards}

\subsection{\S~75030. Minimum Staffing Requirements}

\begin{enumerate}[label=(\alph*)]
\item An authorized Physical AI oncology clinical trial site operating on a
  24-hour schedule shall maintain the following minimum staffing:
  \begin{enumerate}[label=(\arabic*)]
  \item Two to three site safety officers per shift, with three shifts covering
    the 24-hour operating period.
  \item One on-call emergency physician available within 15 minutes at all
    times.
  \item One on-call pharmacist available within 15 minutes at all times.
  \item One facility maintenance technician on-site per shift.
  \item One cybersecurity operations specialist on-call at all times.
  \end{enumerate}

\item Total full-time equivalent staffing shall be a minimum of 8 and a maximum
  of 10, reflecting the Physical AI model in which more patients are treated
  with more robots and fewer workers~\cite{sitespec2026}.

\item The site shall maintain a staffing contingency plan for situations
  requiring additional human personnel, such as multiple simultaneous
  adverse events or extended system outages.
\end{enumerate}

\subsection{\S~75031. Training and Certification}

\begin{enumerate}[label=(\alph*)]
\item Site safety officers shall complete a department-approved training program
  of not less than 160 hours covering the following:
  \begin{enumerate}[label=(\arabic*)]
  \item Physical AI system operation and capabilities for all 10 robot
    categories deployed at the site.
  \item Emergency stop procedures and manual override protocols.
  \item Adverse event recognition, documentation, and reporting.
  \item Patient interaction protocols, including patient rights under AB 2847.
  \item Cybersecurity awareness and incident response.
  \item Regulatory compliance including ICH E6(R3), 21 CFR Part 50, and
    21 CFR Part 312 as adapted for Physical AI~\cite{iche6r3adapt,cfr50adapt,cfr312adapt}.
  \end{enumerate}

\item Certification shall be valid for two years, with 40 hours of continuing
  education required for renewal.

\item The department shall develop and maintain the training curriculum in
  consultation with Physical AI system manufacturers, oncology professional
  organizations, and patient advocacy groups.
\end{enumerate}

\section{Inspection and Enforcement}

\subsection{\S~75040. Routine Inspections}

\begin{enumerate}[label=(\alph*)]
\item The department shall conduct routine inspections of authorized sites at
  least quarterly, without prior notice.

\item Inspections shall assess compliance with authorization conditions,
  Physical AI system registration and maintenance, staffing requirements,
  patient rights protections, adverse event reporting, data protection, and
  facility condition.

\item The department shall publish inspection reports on its website within 30
  days of the inspection.
\end{enumerate}

\subsection{\S~75041. Enforcement Actions}

\begin{enumerate}[label=(\alph*)]
\item The department may take the following enforcement actions for
  noncompliance:
  \begin{enumerate}[label=(\arabic*)]
  \item Written notice of deficiency with required corrective action plan.
  \item Administrative penalty of up to ten thousand dollars (\$10,000) per
    day per violation.
  \item Suspension of authorization to use specific Physical AI system types.
  \item Suspension of site authorization.
  \item Revocation of site authorization.
  \end{enumerate}

\item Before suspending or revoking authorization, the department shall provide
  written notice and an opportunity for hearing in accordance with the
  Administrative Procedure Act (Government Code \S~11500 et seq.).

\item In cases of imminent danger to patient safety, the department may issue
  an immediate suspension order, with a hearing to be held within 15 days.
\end{enumerate}

\section{Reporting Requirements}

\subsection{\S~75050. Monthly Reports}

\begin{enumerate}[label=(\alph*)]
\item Authorized sites shall submit monthly reports to the department including:
  \begin{enumerate}[label=(\arabic*)]
  \item Total patients served and procedures performed, disaggregated by
    cancer type and Physical AI system type.
  \item Current PSL scores for all robot types and cumulative site PSL score.
  \item All adverse events with grade, resolution, and root cause analysis.
  \item Physical AI system utilization rates and maintenance records.
  \item Staffing levels and any departures from minimum requirements.
  \end{enumerate}
\end{enumerate}

\subsection{\S~75051. Annual Reports}

\begin{enumerate}[label=(\alph*)]
\item Authorized sites shall submit annual reports including all elements
  required under the adapted 21 CFR Part 312 \S~312.33, plus the following
  state-specific elements~\cite{cfr312adapt}:
  \begin{enumerate}[label=(\arabic*)]
  \item Year-over-year comparison of patient outcomes and throughput.
  \item Updated USL scores for all Physical AI systems with trend analysis.
  \item Community benefits achieved, including local workforce development.
  \item Patient satisfaction data and complaint resolution statistics.
  \item Comparison of site performance to the baseline established by the
    24-hour simulation (168 patients, 99.7 percent uptime, 8-minute average
    wait, 4.7 out of 5.0 satisfaction)~\cite{newtrial2026}.
  \end{enumerate}
\end{enumerate}


% ============================================================================
% DOCUMENT 6: FDA Physical AI Oncology Trial Site National Compliance Guide
% ============================================================================

\part{FDA Physical AI Oncology Trial Site National Compliance Guide}
\setcounter{section}{0}

\section{Scope and Purpose}

\subsection{Scope}

This compliance guide identifies the federal regulations applicable to
Physical AI oncology clinical trial sites and specifies where corrections,
additions, and adaptations should be applied to existing regulatory text.
Because a comprehensive amendment of each federal regulation is beyond the
scope of a single document, this guide uses a reference-and-correction approach:
for each applicable regulation, the specific sections requiring updates are
identified, and the nature of the required changes is described.

\subsection{Purpose}

The purpose is to provide Physical AI oncology clinical trial site operators,
sponsors, and regulators with a consolidated roadmap for federal compliance.
This guide is informed by three detailed regulatory adaptations already
completed:

\begin{enumerate}[label=(\arabic*)]
\item Adaptation of ICH E6(R3) for Physical AI oncology clinical trials
  (1,301 lines, 50+ defined terms), covering good clinical practice,
  investigator responsibilities, sponsor responsibilities, and data
  governance~\cite{iche6r3adapt}.
\item Adaptation of 21 CFR Part 50 for Physical AI oncology clinical trials
  (748 lines, 18 Physical AI-specific definitions), covering human subjects
  protection and informed consent~\cite{cfr50adapt}.
\item Adaptation of 21 CFR Part 312 for Physical AI oncology clinical trials
  (1,700 lines, 22 Physical AI definitions), covering IND applications,
  graduated autonomy, safety reporting, and
  decommissioning~\cite{cfr312adapt}.
\end{enumerate}

\subsection{Evidence Base}

The evidence for a Physical AI oncology trial transition is overwhelmingly
based on the most recent simulation~\cite{newtrial2026}, which demonstrated:

\begin{itemize}
\item 168 unique patients across 15 cancer types served in 24 hours by 29
  robot instances of 10 types.
\item Patient throughput 10 times greater than traditional trials (150 to 180
  patients per day versus 5 to 15).
\item Per-patient cost reduction of 75 to 85 percent.
\item Patient retention improvement from 70 to 80 percent (traditional) to
  85 to 95 percent (Physical AI).
\item 99.7 percent facility uptime across 24 continuous hours.
\item Average patient wait time of 8 minutes.
\item 7 adverse events (all Grade 1 or 2, all resolved same hour).
\item Patient satisfaction of 4.7 out of 5.0.
\item Cumulative site PSL of 63.4 to 64.4 (Advanced Site band).
\item Both minute-level clinical detail and 24-hour operational continuity
  handled correctly by AI systems.
\item Staffing reduced from 80 to 120 FTE to 8 to 10 FTE.
\item All clinical procedures, monitoring, documentation, and regulatory
  compliance performed autonomously.
\item Seven surgical procedures with 100 percent R0 resection rate.
\item Peak utilization of 72.4 percent in hour 09 with 28 concurrent patients.
\end{itemize}

\section{21 CFR Part 50 -- Protection of Human Subjects}

\subsection{Sections Requiring Update}

The adapted 21 CFR Part 50 (DOI: 10.5281/zenodo.19040707) provides the
comprehensive text~\cite{cfr50adapt}. The following sections of the original
regulation require corrections:

\begin{enumerate}[label=(\alph*)]
\item \textbf{\S~50.1 (Scope):} Add Physical AI system categories and
  applicability criteria. Specify that the regulation applies to all clinical
  investigations where Physical AI systems perform, assist, or monitor
  procedures. Add minimum USL thresholds by robot type.

\item \textbf{\S~50.3 (Definitions):} Add 18 Physical AI-specific definitions
  including robot agent, USL, digital twin, simulation framework, MCP,
  hash-chained audit trail, HIPAA Safe Harbor de-identification, PCCP,
  federated learning, pre-procedure safety matrix, and conformance levels.

\item \textbf{\S~50.20 (General Requirements for Informed Consent):} Extend
  to require disclosure of robotic system role, autonomy mode, USL score in
  plain language, safety record, emergency stop provisions, digital twin usage,
  and MCP data processing pipeline.

\item \textbf{\S~50.25 (Elements of Informed Consent):} Add Physical AI
  elements as additional required elements, including all 10 robot categories
  and their patient-facing instruction formats~\cite{patientinstructions2026}.

\item \textbf{\S~50.50 through \S~50.56 (Pediatric Protections):} Add
  age-appropriate robotic demonstration requirements, parental consent
  provisions for Physical AI procedures, and advocate appointments.
\end{enumerate}

\subsection{Implementation Notes}

Sites should apply the full adapted text from the Physical AI 21 CFR Part 50
document~\cite{cfr50adapt}. Where the adapted text extends the original
regulation, the extended provisions are additive and do not replace the
original human subjects protection requirements.

\section{21 CFR Part 312 -- Investigational New Drug Application}

\subsection{Sections Requiring Update}

The adapted 21 CFR Part 312 (DOI: 10.5281/zenodo.19057628) provides the
comprehensive text~\cite{cfr312adapt}. Key sections requiring corrections:

\begin{enumerate}[label=(\alph*)]
\item \textbf{\S~312.1-312.2 (Scope):} Extend to cover clinical investigations
  where Physical AI systems perform, assist, or monitor investigational drug
  administration. Define when IND coverage is required based on Physical AI
  risk introduction.

\item \textbf{\S~312.3 (Definitions):} Add 22 Physical AI definitions
  including robot agents, USL scoring (9 platforms evaluated), digital twins,
  simulation frameworks, robot capability profiles, task orders, emergency
  stop, human oversight, MCP, and VLA models.

\item \textbf{\S~312.21 (Phases of Investigation):} Add Phase 0 simulation
  validation requirements. Establish graduated autonomy: Phase 1 teleoperation
  only, Phase 2 graduated autonomy with safety data review, Phase 3 full
  autonomy range with accumulated evidence~\cite{cfr312adapt}.

\item \textbf{\S~312.23 (IND Content):} Add Physical AI System Description
  requirements: system identification, USL score, role in drug administration,
  hardware and software specifications, PCCP, Phase 0 results, digital twin
  specifications, safety architecture, cybersecurity assessment, MCP
  integration, human-robot interaction protocols, and maintenance procedures.

\item \textbf{\S~312.30 (Protocol Amendments):} Define Physical AI-specific
  amendments: new robot type, autonomy level change, significant AI model
  update, simulation framework change, human oversight change, digital twin
  methodology change, and MCP architecture change.

\item \textbf{\S~312.32 (Safety Reporting):} Add seven Physical AI adverse
  event categories: serious Physical AI adverse events, Physical AI-drug
  interactions, cybersecurity incidents (7-day or 15-day reporting), AI model
  performance degradation, sim-to-real divergence (trajectory above 4mm or
  force above 1.0N), digital twin discrepancy, and Physical AI Safety Review
  Committee findings.

\item \textbf{\S~312.33 (Annual Reports):} Add Physical AI elements: USL
  score updates, AI model updates, simulation results, digital twin metrics,
  cybersecurity events, and safety committee findings.

\item \textbf{\S~312.38 (Withdrawal of IND):} Add decommissioning
  requirements: safe deactivation, log preservation, secure patient data
  erasure, 30-day decommissioning timeline, 60-day technical
  analysis~\cite{cfr312adapt}.

\item \textbf{\S~312.40 (Requirements for Use):} Add Physical AI readiness
  requirements: system description submitted, USL threshold met, pre-procedure
  safety matrix verified, training completed, MCP infrastructure configured.

\item \textbf{\S~312.42 (Clinical Holds):} Add eight Physical AI-specific
  grounds: robotic safety failures, AI model degradation, simulation-reality
  divergence (above 2x tolerance), cybersecurity compromise, USL score
  decline, inadequate system description, digital twin failure, and
  inadequate human oversight.
\end{enumerate}

\section{ICH E6(R3) -- Good Clinical Practice}

\subsection{Sections Requiring Update}

The adapted ICH E6(R3) (DOI: 10.5281/zenodo.18973368) provides the
comprehensive text~\cite{iche6r3adapt}. Key sections requiring corrections:

\begin{enumerate}[label=(\alph*)]
\item \textbf{Section 1 (Principles):} Extend foundational principles to
  encompass robotic surgical systems, AI algorithms, and digital twins.
  Add Physical AI system classification (seven robot categories) and the USL
  scoring framework (v1.4.0 through v1.8.0)~\cite{usl2026}.

\item \textbf{Section 2 (Investigator Responsibilities):} Add robotic system
  oversight, AI model monitoring, digital twin validation, and cybersecurity
  incident management to investigator duties.

\item \textbf{Section 3 (Sponsor Responsibilities):} Add quality management
  for Physical AI components: robotic system selection, AI lifecycle
  management, simulation validation, digital twin deployment, and federated
  multi-site coordination.

\item \textbf{Section 4 (Data Governance):} Add full lifecycle governance for
  Physical AI data: robotic telemetry, AI decision logs, digital twin state
  data. Mandate hash-chained audit trails (SHA-256) and HIPAA Safe Harbor
  de-identification. Address five AI categories: generative AI (VLA, diffusion
  policies), agentic AI (LLM-based), reinforcement learning, self-supervised
  learning, and supervised learning.

\item \textbf{Appendix A (Physical AI System Documentation):} Add requirements
  analogous to Investigator's Brochure for Physical AI system descriptions.

\item \textbf{Appendix B (Protocol Templates):} Add Physical AI-specific
  protocol elements including simulation validation, robot deployment plans,
  and graduated autonomy phases.

\item \textbf{Glossary:} Add 50+ terms defined including agentic AI,
  calibration, cobots, digital twins, federated learning, sim-to-real
  transfer, and USL framework.
\end{enumerate}

\section{21 CFR Part 11 -- Electronic Records}

\subsection{Sections Requiring Update}

\begin{enumerate}[label=(\alph*)]
\item \textbf{\S~11.10 (Controls for Closed Systems):} Specify that Physical
  AI system telemetry, AI decision logs, and digital twin records are
  electronic records subject to Part 11 requirements. Hash-chained audit
  trails (SHA-256) satisfy the audit trail requirement when cryptographic
  chain integrity is verified quarterly~\cite{iche6r3adapt}.

\item \textbf{\S~11.30 (Controls for Open Systems):} Apply to federated
  learning data exchanges between Physical AI systems at different sites,
  requiring encryption and digital signature verification.

\item \textbf{\S~11.50 (Signature Manifestations):} Extend to include
  automated system signatures from Physical AI systems that generate,
  modify, or transmit electronic records during clinical procedures.
\end{enumerate}

\section{21 CFR Part 820 -- Quality System Regulation}

\subsection{Sections Requiring Update}

\begin{enumerate}[label=(\alph*)]
\item \textbf{\S~820.30 (Design Controls):} Apply to Physical AI systems
  classified as medical devices, including Phase 0 simulation validation as
  part of design verification and validation.

\item \textbf{\S~820.70 (Production and Process Controls):} Apply to the
  deployment and configuration of Physical AI systems at trial sites,
  including calibration, maintenance, and software update processes.

\item \textbf{\S~820.90 (Nonconforming Product):} Apply to Physical AI
  systems that fall below minimum USL thresholds or exhibit performance
  degradation, with corrective action timelines aligned with the adapted
  21 CFR Part 312 safety reporting requirements~\cite{cfr312adapt}.

\item \textbf{\S~820.198 (Complaint Files):} Extend to include Physical AI
  system malfunction reports, patient complaints involving robotic systems,
  and AI model performance complaints.
\end{enumerate}

\section{Additional Federal Requirements}

\subsection{HIPAA (45 CFR Parts 160 and 164)}

Physical AI oncology clinical trial sites shall comply with HIPAA Privacy,
Security, and Breach Notification Rules with additional attention to the
following Physical AI-specific considerations:

\begin{enumerate}[label=(\alph*)]
\item Robotic telemetry data that can be linked to specific patients
  constitutes protected health information (PHI) under HIPAA.
\item Digital twin data constitutes PHI when it contains or is derived from
  individually identifiable health information.
\item MCP-PAI server architectures shall implement the HIPAA Security Rule
  administrative, physical, and technical safeguards~\cite{iche6r3adapt}.
\item The five MCP-PAI servers (authorization, FHIR, DICOM, ledger,
  provenance) shall each maintain individual HIPAA compliance
  documentation.
\end{enumerate}

\subsection{Nuclear Regulatory Commission (10 CFR Part 35)}

Physical AI oncology clinical trial sites operating radiotherapy systems
shall comply with NRC radiation safety requirements, with the following
Physical AI-specific considerations:

\begin{enumerate}[label=(\alph*)]
\item Radiotherapy positioning robots and motion-tracking robots shall be
  included in the site's radiation safety program as radiation-producing
  equipment accessories.
\item The radiation safety officer shall be trained in Physical AI system
  radiation safety implications.
\item Radiation dose calculations shall account for the precision of Physical
  AI positioning systems, which demonstrated submillimeter accuracy and 94.2
  percent beam gating efficiency in the
  simulation~\cite{newtrial2026}.
\end{enumerate}

\section{Cross-Reference Matrix}

The following table maps Physical AI system categories to the primary federal
regulations that apply to each:

\begin{center}
\footnotesize
\begin{tabular}{lccccc}
\toprule
\textbf{System Type} & \textbf{Part 50} & \textbf{Part 312} &
\textbf{E6(R3)} & \textbf{Part 11} & \textbf{Part 820} \\
\midrule
Surgical robots       & Yes & Yes & Yes & Yes & Yes \\
Cobots                & Yes & Yes & Yes & Yes & Yes \\
RT positioning        & Yes & Yes & Yes & Yes & Yes \\
Needle-placement      & Yes & Yes & Yes & Yes & Yes \\
Social companion      & Yes & --  & Yes & Yes & -- \\
Humanoids             & Yes & --  & Yes & Yes & -- \\
RT motion-tracking    & Yes & Yes & Yes & Yes & Yes \\
Imaging assistants    & Yes & Yes & Yes & Yes & Yes \\
Steerable needles     & Yes & Yes & Yes & Yes & Yes \\
Rehab exoskeletons    & Yes & --  & Yes & Yes & Yes \\
\bottomrule
\end{tabular}
\end{center}


% ============================================================================
% DOCUMENT 7: Physical AI Oncology Clinical Trial Facility Building Code Standards
% ============================================================================

\part{Physical AI Oncology Clinical Trial Facility Building Code Standards}
\setcounter{section}{0}

\section{General Requirements}

\subsection{Scope}

These building code standards supplement the California Building Standards Code
(Title 24, Parts 1 through 12) with additional requirements specific to Physical
AI oncology clinical trial facilities. These standards apply to new construction
of facilities designed to house 10 categories of Physical AI robotic systems
operating autonomously on a 24-hour, 7-day-per-week schedule serving up to 180
oncology patients per day~\cite{sitespec2026}.

\subsection{Facility Size and Configuration}

\begin{enumerate}[label=(\alph*)]
\item The minimum gross floor area shall be 85,000 square feet, with a maximum
  of 100,000 square feet, distributed across the clinical wings, support areas,
  and infrastructure spaces specified in these standards~\cite{sitespec2026}.

\item The facility shall accommodate nine clinical wings corresponding to the
  following Physical AI system deployment zones:
  \begin{enumerate}[label=(\arabic*)]
  \item Three surgical suites for surgical robots (SURG-01, SURG-02, SURG-03).
  \item Four cobot stations for collaborative biopsy robots.
  \item Three radiotherapy vaults for RT positioning robots.
  \item Two needle-placement suites for robotic needle-placement systems.
  \item Five social companion stations for pediatric and adult companion
    interactions.
  \item Three humanoid stations for rehabilitation engagement.
  \item Three RT motion-tracking vaults (co-located with RT positioning).
  \item Four imaging bays for imaging assistant robots.
  \item Two steerable needle suites.
  \item One rehabilitation exoskeleton bay with two active stations and one
    reserved station.
  \end{enumerate}

\item Support areas shall include a robot-accessible pharmacy, server room,
  maintenance bay, and decontamination area~\cite{newtrial2026}.
\end{enumerate}

\section{Structural Requirements}

\subsection{Floor Load Capacity}

\begin{enumerate}[label=(\alph*)]
\item Surgical suites shall be designed for a minimum floor live load of 150
  pounds per square foot to support surgical robot systems (including the
  da Vinci Xi surgical system at approximately 1,800 pounds), instrument carts,
  patient tables, and support equipment~\cite{patientinstructions2026}.

\item Radiotherapy vaults shall be designed for a minimum floor live load of
  200 pounds per square foot to support linear accelerator systems,
  radiotherapy positioning robots, and radiation shielding components.

\item Imaging bays shall be designed for a minimum floor live load of 125
  pounds per square foot to support imaging equipment, robotic ultrasound
  systems, and CT scanners used for needle-placement guidance.

\item Rehabilitation exoskeleton areas shall include floor-mounted parallel
  bars with anchorage rated for 500 pounds lateral force, supporting patients
  during exoskeleton-assisted gait
  training~\cite{patientinstructions2026}.

\item Server rooms shall be designed for a minimum floor live load of 250
  pounds per square foot to support server rack installations for AI
  computation infrastructure.

\item General clinical circulation areas shall maintain a minimum floor live
  load of 100 pounds per square foot, accommodating hospital beds,
  wheelchairs, and mobile robot systems simultaneously.
\end{enumerate}

\subsection{Vibration Control}

\begin{enumerate}[label=(\alph*)]
\item Surgical suites and needle-placement suites shall meet VC-E vibration
  criteria (peak velocity 3.1 micrometers per second) to support the
  submillimeter precision required by surgical robots and needle-placement
  systems~\cite{newtrial2026}.

\item Imaging bays shall meet VC-D vibration criteria (peak velocity 6.3
  micrometers per second) to support diagnostic imaging quality.

\item Radiotherapy vaults shall meet VC-C vibration criteria (peak velocity
  12.5 micrometers per second) to support accurate beam alignment.

\item Vibration isolation joints shall separate clinical wings from building
  mechanical systems, loading areas, and parking structures.
\end{enumerate}

\section{Mechanical Systems (HVAC)}

\subsection{Surgical Suite Air Handling}

\begin{enumerate}[label=(\alph*)]
\item Surgical suites shall maintain ISO Class 7 (10,000 particles per cubic
  foot) or better air cleanliness using HEPA filtration with a minimum of 20
  air changes per hour~\cite{sitespec2026}.

\item Surgical suites shall maintain positive pressure of at least 2.5 Pascals
  relative to adjacent spaces.

\item Temperature control shall maintain 68 to 73 degrees Fahrenheit with
  humidity between 30 and 60 percent relative humidity.

\item Independent air handling units shall serve each surgical suite to prevent
  cross-contamination and allow maintenance without affecting other suites.
\end{enumerate}

\subsection{Radiotherapy Vault Ventilation}

\begin{enumerate}[label=(\alph*)]
\item Radiotherapy vaults shall have dedicated exhaust systems capable of
  removing ozone generated by high-energy radiation beams.

\item A minimum of 6 air changes per hour shall be maintained during treatment
  operations, with 10 air changes per hour during purge cycles between patients.

\item Exhaust air from radiotherapy vaults shall not be recirculated to other
  building areas.
\end{enumerate}

\subsection{Server Room Cooling}

\begin{enumerate}[label=(\alph*)]
\item Server rooms supporting AI computation infrastructure shall maintain
  temperature between 64 and 75 degrees Fahrenheit with redundant cooling
  systems providing N+1 capacity~\cite{sitespec2026}.

\item Cooling systems shall reject a minimum heat load of 200 kilowatts,
  calculated from the server room's portion of the facility's total 800 to
  1,200 kilowatt electrical load.

\item Server room cooling shall be on a dedicated HVAC zone independent of
  clinical space climate control.

\item Hot aisle and cold aisle containment shall be used to maximize cooling
  efficiency.
\end{enumerate}

\section{Electrical Systems}

\subsection{Power Supply}

\begin{enumerate}[label=(\alph*)]
\item The facility shall be served by dual utility power feeds from separate
  substations or feeders, providing full redundancy for the 800 to 1,200
  kilowatt total electrical load~\cite{sitespec2026}.

\item Automatic transfer switches shall provide seamless transition between
  utility feeds, with a maximum transfer time of 10 milliseconds for critical
  clinical systems.

\item Emergency generator capacity shall be sufficient to power all critical
  clinical systems (surgical suites, radiotherapy vaults, server rooms,
  life safety systems) for a minimum of 72 hours at full load.

\item Generator fuel storage shall accommodate 96 hours of operation with
  automatic fuel management and monitoring.
\end{enumerate}

\subsection{Uninterruptible Power Supply (UPS)}

\begin{enumerate}[label=(\alph*)]
\item All surgical suites shall be served by dedicated UPS systems providing a
  minimum of 30 minutes of battery backup at full surgical robot operating
  load, sufficient to safely complete or abort any procedure in progress.

\item Server rooms shall be served by UPS systems providing a minimum of 15
  minutes of battery backup, sufficient for orderly shutdown of AI computation
  systems and data preservation.

\item All Physical AI system charging and standby stations shall be served by
  UPS-backed circuits to maintain robot readiness during power transitions.

\item UPS systems shall be tested monthly under load, with results recorded
  in the facility maintenance log.
\end{enumerate}

\subsection{Electrical Distribution}

\begin{enumerate}[label=(\alph*)]
\item Dedicated electrical panels shall serve each clinical wing, with
  sufficient circuit capacity for all Physical AI system types deployed in
  that wing plus 25 percent spare capacity.

\item Isolated ground circuits shall serve all Physical AI systems and medical
  imaging equipment to prevent electromagnetic interference.

\item Power quality monitoring systems shall continuously track voltage, current,
  frequency, harmonics, and transients on all circuits serving Physical AI
  systems, with automated alerting for out-of-specification conditions.
\end{enumerate}

\section{Radiation Shielding}

\begin{enumerate}[label=(\alph*)]
\item Radiotherapy vault shielding shall comply with NCRP Report 151 and
  California Radiologic Health Branch requirements, designed for the maximum
  beam energy and workload planned for the facility.

\item Shielding design shall account for the Physical AI positioning robot's
  range of motion, ensuring that all achievable patient positions result in
  adequate shielding of adjacent spaces.

\item The primary barrier shall be designed for the full workload of the
  radiotherapy system, and secondary barriers shall account for leakage and
  scatter radiation.

\item Maze entrances shall be designed to reduce radiation levels at the maze
  entrance to below 0.02 milliSieverts per hour during beam-on conditions.

\item Shielding calculations shall be reviewed and certified by a qualified
  medical physicist before construction begins.
\end{enumerate}

\section{Technology Infrastructure}

\subsection{Network Architecture}

\begin{enumerate}[label=(\alph*)]
\item The facility shall provide a minimum of 10 gigabits per second internal
  network bandwidth using redundant fiber-optic backbone
  infrastructure~\cite{sitespec2026}.

\item Clinical networks, administrative networks, and Physical AI system control
  networks shall be physically or logically segmented with firewall separation
  between segments.

\item Each clinical wing shall have dedicated network switches with redundant
  uplinks to the core network.

\item Wireless coverage (Wi-Fi 6E or later) shall be provided throughout all
  patient-accessible areas for patient portal access and companion robot
  connectivity.
\end{enumerate}

\subsection{Robot Charging and Docking Infrastructure}

\begin{enumerate}[label=(\alph*)]
\item Each clinical wing shall include dedicated robot charging and docking
  stations rated for the specific power requirements of the Physical AI
  systems deployed in that wing.

\item Charging stations shall include automated connection systems to minimize
  human intervention during shift changes and standby periods.

\item Floor-mounted guide paths or ceiling-mounted positioning references shall
  facilitate autonomous robot navigation between clinical stations, charging
  stations, and maintenance areas.

\item The facility layout shall support the patient flow demonstrated in the
  simulation, where peak utilization reached 21 of 29 robots active
  simultaneously during hour 09, with an average of 8.1 of 29 active across
  24 hours~\cite{newtrial2026}.
\end{enumerate}

\section{Life Safety and Egress}

\begin{enumerate}[label=(\alph*)]
\item The facility shall comply with California Fire Code (Title 24, Part 9)
  and NFPA 101 Life Safety Code requirements for health care occupancies.

\item Emergency egress paths shall be designed to accommodate patients on
  hospital beds and in wheelchairs, with corridor widths of at least 8 feet
  in primary clinical circulation areas.

\item Physical AI systems shall be designed to move to designated safe positions
  during fire alarm activation, clearing egress paths and de-energizing
  non-essential systems.

\item Fire suppression systems in surgical suites shall use clean agent (FM-200
  or Novec 1230) to protect sensitive robotic equipment while maintaining
  patient safety.

\item Fire suppression in server rooms shall use clean agent systems with
  pre-action activation to prevent accidental discharge damage to computing
  equipment.
\end{enumerate}


% ============================================================================
% DOCUMENT 8: Physical AI Oncology Clinical Trial Site Premises Code
% ============================================================================

\part{Physical AI Oncology Clinical Trial Site Premises Code}
\setcounter{section}{0}

\section{Scope and Application}

These premises standards govern the physical site, grounds, security systems,
environmental controls, and operational conditions for Physical AI oncology
clinical trial facilities authorized under California law. The standards address
the unique premises requirements created by 24-hour autonomous robot operation,
patient-centric on-demand scheduling, and the deployment of 29 Physical AI
system instances across a 85,000 to 100,000 square-foot clinical
facility~\cite{sitespec2026}.

The premises code complements the building code standards (structural,
mechanical, electrical, technology) by addressing how the facility and its
grounds are maintained, secured, and operated during continuous clinical
operations. The 24-hour simulation demonstrated that sustained operations
require premises management that handles both minute-level details (individual
patient movements, robot state transitions, adverse event responses) and
facility-wide patterns across a full 24-hour cycle~\cite{newtrial2026}.

\section{Site Security}

\subsection{Physical Security Systems}

\begin{enumerate}[label=(\alph*)]
\item The premises shall be equipped with a comprehensive security system
  including the following components:
  \begin{enumerate}[label=(\arabic*)]
  \item Video surveillance covering all exterior entrances, parking areas,
    patient drop-off zones, lobby and reception areas, clinical wing corridors,
    and robot maintenance areas.
  \item Electronic access control on all exterior doors, all clinical wing
    entries, server rooms, pharmacy, robot maintenance bay, and any area
    housing controlled substances or sensitive equipment.
  \item Intrusion detection systems on all exterior doors, ground-floor windows,
    and roof access points.
  \item Duress alarm stations at the reception desk, each shift supervisor
    station, and the security monitoring room.
  \end{enumerate}

\item Security systems shall operate continuously during the 24-hour operating
  period without interruption.

\item Security camera recordings shall be retained for a minimum of 90 days
  and made available to the department and law enforcement upon request.
\end{enumerate}

\subsection{Access Control Zones}

The premises shall be divided into the following access control zones:

\begin{enumerate}[label=(\alph*)]
\item \textbf{Public Zone:} Lobby, reception area, patient waiting areas,
  restrooms, and the covered patient drop-off canopy. Accessible to patients,
  visitors, and the public during operating hours.

\item \textbf{Patient Zone:} Clinical wings, examination rooms, procedure
  suites, recovery areas, and patient circulation corridors. Accessible to
  enrolled patients with valid credentials and accompanying persons.

\item \textbf{Clinical Zone:} Areas where Physical AI systems perform clinical
  procedures, including surgical suites, RT vaults, needle-placement suites,
  and steerable needle suites. Restricted to authorized clinical personnel,
  site safety officers, and patients undergoing procedures.

\item \textbf{Technical Zone:} Server rooms, robot maintenance bay,
  decontamination area, and equipment storage. Restricted to authorized
  technical personnel and site safety officers.

\item \textbf{Administrative Zone:} Staff offices, break rooms, training rooms,
  and supply storage. Restricted to site personnel.
\end{enumerate}

\subsection{Cybersecurity of Physical Premises}

\begin{enumerate}[label=(\alph*)]
\item Physical AI system control networks shall be accessible only from within
  the Technical Zone, with no wireless access points connected to robot control
  systems.

\item The MCP-PAI server infrastructure shall be housed in the Technical Zone
  server room with physical access limited to authorized cybersecurity and
  system administration personnel~\cite{iche6r3adapt}.

\item Network infrastructure components (switches, routers, fiber distribution
  panels) shall be housed in locked enclosures within the Technical Zone.

\item Visitor and patient wireless networks shall be physically isolated from
  clinical and Physical AI system networks.
\end{enumerate}

\section{Patient Access and Flow}

\subsection{Arrival and Reception}

\begin{enumerate}[label=(\alph*)]
\item The main patient entrance shall be accessible 24 hours per day, 7 days
  per week, consistent with the on-demand scheduling model that demonstrated
  patient arrivals across all 24 hours of the simulation, including 15
  arrivals during the peak hour (hour 09) and continued arrivals during
  overnight hours~\cite{newtrial2026}.

\item A covered patient drop-off zone shall accommodate a minimum of three
  vehicles simultaneously, with a covered walkway extending at least 20 feet
  to the building entrance.

\item The reception area shall include automated check-in kiosks that verify
  patient identity, confirm scheduled procedures, and direct patients to the
  appropriate clinical wing.

\item The reception area shall include seating for a minimum of 20 patients and
  accompanying persons, reflecting the simulation's peak concurrent on-site
  population of 28 patients~\cite{newtrial2026}.
\end{enumerate}

\subsection{Internal Patient Circulation}

\begin{enumerate}[label=(\alph*)]
\item Primary patient circulation corridors shall be a minimum of 8 feet wide
  to accommodate hospital beds, wheelchairs, rehabilitation exoskeletons, and
  mobile Physical AI systems simultaneously.

\item Wayfinding shall use a combination of visual signage, floor markings,
  digital displays, and optional auditory guidance compatible with the 10
  Physical AI system categories that patients may
  encounter~\cite{patientinstructions2026}.

\item Patient flow through the facility shall follow the patterns validated in
  the simulation, with separate intake, procedure, and recovery pathways to
  minimize cross-traffic and wait times~\cite{newtrial2026}.

\item Recovery areas adjacent to each procedure zone shall accommodate patients
  during the post-procedure observation periods specified in the patient robot
  instructions (ranging from 2 minutes for imaging to 60 minutes for steerable
  needle and surgical procedures)~\cite{patientinstructions2026}.
\end{enumerate}

\section{Environmental Standards}

\subsection{Exterior Grounds}

\begin{enumerate}[label=(\alph*)]
\item The site shall maintain landscaping and hardscaping consistent with the
  character of the surrounding neighborhood and the San Francisco Planning
  Code requirements for the applicable zoning district.

\item Exterior lighting shall provide a minimum of 5 foot-candles at all
  pedestrian pathways, parking areas, and building entrances, ensuring patient
  safety during the 24-hour operating period.

\item A minimum of 10 percent of the total lot area shall be dedicated to
  landscaped open space, which may include patient garden areas accessible from
  the Public Zone.

\item Signage shall clearly identify the facility as a Physical AI oncology
  clinical trial site, with the authorization certificate number displayed at
  the main entrance.
\end{enumerate}

\subsection{Interior Environment}

\begin{enumerate}[label=(\alph*)]
\item Patient waiting areas and recovery spaces shall maintain interior
  temperatures between 68 and 74 degrees Fahrenheit with ambient noise levels
  not exceeding 45 dBA.

\item Lighting in patient areas shall be adjustable between 200 and 500 lux to
  accommodate different times of day and patient comfort needs during 24-hour
  operations.

\item Clinical procedure areas shall maintain lighting appropriate to the
  specific Physical AI system requirements, ranging from high-intensity surgical
  lighting in surgical suites to reduced ambient lighting in imaging bays and
  RT vaults.

\item Air quality in all patient-accessible areas shall be maintained at a
  minimum of MERV 13 filtration with continuous monitoring of CO2, temperature,
  and humidity.
\end{enumerate}

\section{Robot Operational Zones}

\subsection{Robot Movement Corridors}

\begin{enumerate}[label=(\alph*)]
\item Designated robot movement corridors shall connect clinical wings to the
  robot maintenance bay, charging stations, and decontamination area.

\item Robot corridors shall be separated from primary patient circulation
  pathways where possible, using time-of-day scheduling where physical
  separation is not feasible.

\item Floor surfaces in robot corridors shall be smooth, level (maximum slope
  1:48), and free of thresholds or transitions that could impede autonomous
  robot navigation.

\item Robot corridors shall be equipped with proximity sensors and visual
  warning indicators (flashing lights) to alert human occupants when robots
  are in transit.
\end{enumerate}

\subsection{Robot Maintenance Bay}

\begin{enumerate}[label=(\alph*)]
\item The robot maintenance bay shall accommodate a minimum of three Physical
  AI systems simultaneously for scheduled maintenance, calibration, or repair.

\item The bay shall include dedicated workstations with electrical, compressed
  air, and network connections appropriate for each Physical AI system
  category.

\item Spare parts storage shall be maintained for critical components of all
  10 Physical AI system types, with inventory levels sufficient to support
  the maintenance schedule demonstrated in the simulation (where SURG-01
  maintenance was conducted from 22:15 to 01:59 without disrupting clinical
  operations)~\cite{newtrial2026}.

\item The maintenance bay shall be acoustically isolated from patient areas to
  allow maintenance activities during nighttime hours without disturbing
  patients in adjacent recovery areas.
\end{enumerate}

\section{Emergency Access and Evacuation}

\begin{enumerate}[label=(\alph*)]
\item Emergency vehicle access shall be provided from at least two directions,
  with a minimum 26-foot-wide fire access lane along the full length of the
  building.

\item A dedicated ambulance bay shall accommodate a minimum of two ambulances
  simultaneously, separate from the patient drop-off zone.

\item Evacuation plans shall account for the specific needs of oncology patients
  who may be connected to treatment equipment, recovering from procedures, or
  using rehabilitation exoskeletons at the time of evacuation.

\item Physical AI systems shall execute automated safe-shutdown and position
  procedures upon receiving evacuation signals, including retracting surgical
  arms, lowering treatment couches, and moving to designated safe positions
  to clear evacuation routes~\cite{patientinstructions2026}.

\item Evacuation drills shall be conducted quarterly, including scenarios
  specific to Physical AI system emergencies such as robot control system
  failure, surgical suite power loss during a procedure, and cybersecurity
  incidents requiring physical network isolation.
\end{enumerate}

\section{Waste Management}

\begin{enumerate}[label=(\alph*)]
\item Medical waste generated by Physical AI-assisted oncology procedures shall
  be managed in accordance with the California Medical Waste Management Act
  (Health and Safety Code \S~117600 et seq.).

\item Radioactive waste from radiotherapy operations shall be managed in
  accordance with applicable NRC and California Radiologic Health Branch
  requirements.

\item Robot decontamination waste shall be classified and managed based on the
  type of contamination (biological, chemical, or radioactive).

\item Dedicated waste staging areas shall be located in the Technical Zone with
  separate containers for biohazardous, radioactive, pharmaceutical, and
  general waste.

\item Waste collection shall be scheduled to avoid conflicts with patient
  arrival and departure patterns, using the patient flow data from the
  simulation to identify optimal collection windows~\cite{newtrial2026}.
\end{enumerate}


% ============================================================================
% DOCUMENT 9: Physical AI Oncology Clinical Trial Site Parking and Patient Transportation Standards
% ============================================================================

\part{Physical AI Oncology Clinical Trial Site Parking and Patient Transportation Standards}
\setcounter{section}{0}

\section{Scope and Purpose}

These standards govern the design, construction, and operation of dedicated
parking facilities and patient transportation infrastructure for Physical AI
oncology clinical trial sites authorized in California. The parking facility
is a new-construction structure purpose-built to serve the trial site, designed
to accommodate the unique patient flow patterns of a 24-hour, on-demand
oncology clinical trial.

The 24-hour simulation demonstrated arrival patterns ranging from 2 patients
during overnight hours (hours 02 through 04) to 15 patients during peak hour
(hour 09), with a total of 168 unique patients in 24 hours and a peak
concurrent on-site population of 28 patients~\cite{newtrial2026}. The parking
facility must accommodate these variable demand patterns while maintaining the
8-minute average wait time and the 24/7 accessibility that define the Physical
AI trial format~\cite{formatcomparison2026}.

\section{Parking Facility Design}

\subsection{Capacity Requirements}

\begin{enumerate}[label=(\alph*)]
\item The parking facility shall provide a minimum of 120 parking spaces,
  calculated as follows:
  \begin{enumerate}[label=(\arabic*)]
  \item 80 patient and visitor spaces, based on the peak concurrent on-site
    population of 28 patients plus accompanying persons and a turnover factor
    reflecting variable procedure durations (15 minutes for imaging to 180
    minutes for surgery)~\cite{newtrial2026,patientinstructions2026}.
  \item 15 staff spaces for the 8 to 10 full-time equivalent staff plus
    on-call personnel across three shifts~\cite{sitespec2026}.
  \item 10 accessible spaces (8.3 percent of total), exceeding the minimum
    ADA and California Building Code Chapter 11B requirements for medical
    facilities.
  \item 5 short-term patient drop-off and pick-up spaces adjacent to the
    building entrance.
  \item 10 spaces designated for ride-share and medical transport vehicle
    staging.
  \end{enumerate}

\item The facility shall accommodate future expansion of up to 50 percent
  additional capacity through structural design that supports the addition
  of one parking level.
\end{enumerate}

\subsection{Structure and Configuration}

\begin{enumerate}[label=(\alph*)]
\item The parking facility shall be a new-construction structure, which may be
  either a surface lot or a multi-level structure (not to exceed three levels
  above grade), depending on the site constraints and zoning requirements.

\item If a multi-level structure is constructed, all patient parking levels
  shall be served by a minimum of two elevators and one ramp system meeting
  ADA accessibility standards.

\item The facility shall provide a covered pedestrian connection from the
  parking structure to the clinical building entrance, with a maximum walking
  distance of 300 feet from the most distant parking space to the building
  entrance.

\item Clear ceiling height shall be a minimum of 8 feet 2 inches on all
  levels to accommodate medical transport vehicles and passenger vans.

\item Drive aisles shall be a minimum of 24 feet wide for two-way traffic,
  with parking spaces a minimum of 9 feet wide by 18 feet deep.
\end{enumerate}

\subsection{Accessible Parking}

\begin{enumerate}[label=(\alph*)]
\item Accessible parking spaces shall be located at the level closest to the
  building entrance, with a maximum distance of 100 feet from the space to
  the building entry.

\item A minimum of two van-accessible spaces (8 feet wide with 8-foot access
  aisle) shall be provided.

\item Accessible routes from parking to the building entrance shall be covered,
  illuminated to a minimum of 10 foot-candles, and free of steps, steep grades
  (maximum slope 1:20), or other barriers.

\item Accessible parking signage shall include the international symbol of
  accessibility and shall note that the facility operates 24 hours per day
  for on-demand oncology care.
\end{enumerate}

\section{Patient Drop-Off Zone}

\begin{enumerate}[label=(\alph*)]
\item A dedicated patient drop-off and pick-up zone shall be located immediately
  adjacent to the main building entrance, separated from the parking facility
  access drive.

\item The drop-off zone shall accommodate a minimum of three vehicles
  simultaneously under a covered canopy extending at least 20 feet from the
  building facade.

\item The drop-off zone shall include a level, non-slip surface, curb cuts,
  and seating for patients waiting for transportation.

\item A call button or intercom shall allow patients to request assistance
  from site safety officers for transfer from vehicles to the building
  entrance.

\item The drop-off zone shall be operational 24 hours per day, with adequate
  lighting (minimum 10 foot-candles) and weather protection for all hours of
  operation, consistent with the simulation's demonstrated overnight patient
  arrivals and departures~\cite{newtrial2026}.
\end{enumerate}

\section{Transportation Demand Management}

\subsection{Public Transit Integration}

\begin{enumerate}[label=(\alph*)]
\item The site shall be located within one-quarter mile of a BART station or
  Muni stop, or the site operator shall provide a shuttle service connecting
  the site to the nearest transit hub at intervals not exceeding 30 minutes
  during all operating hours.

\item Real-time transit information displays shall be installed in the patient
  waiting area and the parking facility lobby, showing schedules for BART,
  Muni, and site shuttle services.

\item The site shall participate in the Clipper card program or equivalent
  regional transit payment system to facilitate patient transit use.
\end{enumerate}

\subsection{Ride-Share and Medical Transport}

\begin{enumerate}[label=(\alph*)]
\item The site shall designate a minimum of 10 parking spaces for ride-share
  and medical transport vehicle staging, located near the patient drop-off
  zone but not blocking drop-off lane traffic flow.

\item The site shall coordinate with ride-share services to establish a
  designated pick-up and drop-off geofence that directs drivers to the
  correct facility entrance.

\item For patients who require medical transport following procedures (such as
  those recovering from anesthesia after surgical robot procedures lasting
  90 to 180 minutes), the site shall coordinate discharge scheduling with
  transport providers to minimize patient wait
  times~\cite{patientinstructions2026}.
\end{enumerate}

\subsection{Bicycle and Pedestrian Access}

\begin{enumerate}[label=(\alph*)]
\item The site shall provide a minimum of 20 Class I (enclosed and secured)
  bicycle parking spaces and 10 Class II (outdoor rack) bicycle parking spaces.

\item Bicycle parking shall be located within 100 feet of the building entrance
  and shall include electric bicycle charging capability for at least 4 spaces.

\item Pedestrian access from the public sidewalk to the building entrance shall
  be clearly marked, well-lit (minimum 5 foot-candles), and separated from
  vehicle traffic.
\end{enumerate}

\section{24-Hour Parking Operations}

\subsection{Staffing and Monitoring}

\begin{enumerate}[label=(\alph*)]
\item The parking facility shall be monitored 24 hours per day through video
  surveillance and emergency call stations, consistent with the 24-hour
  operating schedule of the clinical facility.

\item Emergency call stations shall be located at each elevator lobby, at each
  stairwell entrance, and at intervals not exceeding 200 feet in the parking
  area.

\item A parking management system shall track occupancy in real time and display
  available spaces at the facility entrance, directing arriving patients to
  available spaces efficiently.

\item During overnight hours (10:00 PM to 6:00 AM), the parking facility may
  reduce lighting to 75 percent of daytime levels in unoccupied areas while
  maintaining full lighting in occupied areas, pedestrian routes, and the
  patient drop-off zone.
\end{enumerate}

\subsection{Electric Vehicle Support}

\begin{enumerate}[label=(\alph*)]
\item A minimum of 20 percent of parking spaces (24 spaces) shall be equipped
  with Level 2 electric vehicle charging stations (240V, 32A minimum).

\item A minimum of 4 parking spaces shall be equipped with Level 3 DC fast
  charging stations (minimum 50 kW), enabling patients to charge during shorter
  procedures such as imaging (10 to 20 minutes) or companion robot sessions
  (10 to 20 minutes)~\cite{patientinstructions2026}.

\item An additional 20 percent of parking spaces (24 spaces) shall be designed
  as EV-capable, with conduit and panel capacity installed during construction
  to support future charging station installation.

\item EV charging shall be available 24 hours per day, consistent with the
  facility's operating schedule.
\end{enumerate}

\section{Emergency Vehicle Access}

\begin{enumerate}[label=(\alph*)]
\item A dedicated ambulance bay shall be provided separate from both the parking
  facility and the patient drop-off zone, accommodating a minimum of two
  ambulances simultaneously.

\item Fire department access lanes shall be maintained along the full perimeter
  of both the clinical building and the parking facility, with a minimum width
  of 26 feet and a minimum vertical clearance of 13 feet 6 inches.

\item Emergency vehicle access routes shall be clearly marked and kept
  unobstructed at all times, with automatic enforcement through surveillance
  and notification systems.

\item The parking facility design shall not impede emergency vehicle access to
  the clinical building from any direction.
\end{enumerate}

\section{Sustainability}

\begin{enumerate}[label=(\alph*)]
\item The parking facility shall include a stormwater management system
  compliant with the San Francisco Green Infrastructure Ordinance,
  incorporating bioretention areas, permeable pavement, or equivalent
  green infrastructure techniques.

\item Rooftop areas of the parking structure, if applicable, shall incorporate
  either solar photovoltaic panels or a green roof system covering a minimum
  of 50 percent of the roof area.

\item Parking facility lighting shall use LED fixtures with occupancy sensors
  and daylight harvesting to minimize energy consumption during the 24-hour
  operating period.

\item Construction materials shall meet the CalGreen (Title 24, Part 11)
  requirements for recycled content and low-emission materials.
\end{enumerate}


% ============================================================================
% DOCUMENT 10: Physical AI Oncology Clinical Trial Site Activation and Standard Operating Procedures
% ============================================================================

\part{Physical AI Oncology Clinical Trial Site Activation and Standard Operating Procedures}
\setcounter{section}{0}

\section{Site Activation Checklist}

\subsection{Pre-Activation Requirements}

The following requirements shall be completed before the first patient is
enrolled at a Physical AI oncology clinical trial site:

\begin{enumerate}[label=(\alph*)]
\item \textbf{Regulatory Approvals:}
  \begin{enumerate}[label=(\arabic*)]
  \item State authorization certificate issued by CDPH (per SB 1042).
  \item San Francisco conditional use permit and health care facility operating
    license (or equivalent local permits for non-San Francisco sites).
  \item IRB approval for the Physical AI oncology trial protocol.
  \item IND clearance with Physical AI System Description accepted by
    FDA~\cite{cfr312adapt}.
  \item All Physical AI systems registered with CDPH with current USL scores
    meeting minimum thresholds~\cite{usl2026}.
  \end{enumerate}

\item \textbf{Facility Readiness:}
  \begin{enumerate}[label=(\arabic*)]
  \item Building occupancy permit issued.
  \item All Physical AI systems installed, commissioned, and verified.
  \item Phase 0 simulation validation completed for all system types with
    cross-framework consistency documented~\cite{iche6r3adapt}.
  \item Network infrastructure operational with MCP-PAI five-server topology
    configured to the required conformance level.
  \item Power systems tested including UPS and generator transfer.
  \item HVAC systems verified for all clinical zones.
  \item Security systems operational with access control zones configured.
  \end{enumerate}

\item \textbf{Staffing Readiness:}
  \begin{enumerate}[label=(\arabic*)]
  \item All site safety officers trained and CDPH-certified.
  \item On-call emergency physician and pharmacist agreements executed.
  \item Maintenance technician assigned and trained.
  \item Cybersecurity operations specialist on-call roster established.
  \item All staff completed competency assessments for Physical AI systems.
  \end{enumerate}

\item \textbf{Documentation:}
  \begin{enumerate}[label=(\arabic*)]
  \item Standard operating procedures finalized for all 10 robot types.
  \item Patient-facing robot instructions available in English and
    the most common languages spoken in the service
    area~\cite{patientinstructions2026}.
  \item Informed consent documents approved by IRB with Physical AI
    disclosures~\cite{cfr50adapt}.
  \item Emergency response plan posted and distributed.
  \item Adverse event reporting workflows tested.
  \end{enumerate}
\end{enumerate}

\section{Shift Management}

\subsection{Three-Shift Structure}

\begin{enumerate}[label=(\alph*)]
\item The site shall operate on a three-shift schedule covering 24 continuous
  hours:
  \begin{enumerate}[label=(\arabic*)]
  \item Day shift: 06:00 to 14:00.
  \item Evening shift: 14:00 to 22:00.
  \item Night shift: 22:00 to 06:00.
  \end{enumerate}

\item Each shift shall be staffed by a minimum of two site safety officers, with
  a third officer added during peak demand periods (typically 08:00 to 16:00
  based on the simulation's patient arrival patterns showing 15 arrivals
  during hour 09)~\cite{newtrial2026}.

\item Shift handoff procedures shall include a structured briefing covering:
  \begin{enumerate}[label=(\arabic*)]
  \item Current patient census and active procedures.
  \item Physical AI system status (active, standby, maintenance) for all 29
    instances.
  \item Any open adverse events or safety concerns.
  \item Scheduled maintenance windows (such as the SURG-01 22:15 to 01:59
    maintenance demonstrated in the simulation)~\cite{newtrial2026}.
  \item Pending patient arrivals and departures.
  \end{enumerate}
\end{enumerate}

\subsection{Shift Performance Metrics}

Each shift shall track and report the following metrics, consistent with the
hourly data structure validated in the 24-hour simulation:

\begin{enumerate}[label=(\alph*)]
\item Patient arrivals, procedures completed, and departures during the shift.
\item Physical AI system utilization rates per type (the simulation averaged
  8.1 of 29 active, with peak of 21 of 29 active)~\cite{newtrial2026}.
\item Average patient wait time (target: 8 minutes or less).
\item Adverse events with grade, resolution time, and root cause.
\item PSL score changes per robot type per regulatory dimension (maximum
  0.3 change per hour per dimension)~\cite{pslframework2026}.
\item Any emergency stop activations or manual overrides.
\end{enumerate}

\section{Patient Processing Workflow}

\subsection{Arrival to Procedure}

\begin{enumerate}[label=(\alph*)]
\item Patient arrival and check-in via automated kiosk (estimated 2 minutes).
\item Identity verification and informed consent confirmation. Physical AI
  disclosures reviewed for any system changes since enrollment.
\item Pre-procedure safety matrix completed, including verification that all
  assigned Physical AI systems meet USL thresholds and have passed daily
  calibration checks~\cite{cfr50adapt}.
\item Patient escorted or directed to the appropriate clinical wing based on
  scheduled procedure type.
\item Patient-facing robot instructions reviewed with the patient for the
  specific systems involved in their procedure. For example, surgical robot
  patients receive the three-step surgical robot instruction sheet covering
  anesthesia induction, robot-assisted procedure, and post-operative
  recovery~\cite{patientinstructions2026}.
\item Procedure initiated under Physical AI system autonomous control with
  site safety officer oversight.
\end{enumerate}

\subsection{Procedure to Departure}

\begin{enumerate}[label=(\alph*)]
\item Procedure completed; Physical AI system performs automated post-procedure
  protocols (retracting instruments, lowering couches, generating procedure
  reports).
\item Patient transferred to recovery area; recovery duration determined by
  procedure type (2 minutes for imaging to 60 minutes for surgical or
  steerable needle procedures)~\cite{patientinstructions2026}.
\item Post-procedure vital sign monitoring by Physical AI systems until
  discharge criteria met.
\item Discharge documentation generated automatically; patient receives
  procedure summary and follow-up instructions.
\item Patient exits via reception area; transportation coordination if needed.
\end{enumerate}

\section{Physical AI System Daily Operations}

\subsection{Daily System Checks}

\begin{enumerate}[label=(\alph*)]
\item Each Physical AI system shall undergo an automated daily verification
  check before the first patient procedure, including:
  \begin{enumerate}[label=(\arabic*)]
  \item Mechanical calibration verification (positioning accuracy, force
    limits, range of motion).
  \item Software system integrity check (AI model version, parameter
    verification, hash validation).
  \item Network connectivity and MCP-PAI server communication test.
  \item Emergency stop functionality test.
  \item Sensor calibration verification (cameras, force sensors, proximity
    sensors).
  \end{enumerate}

\item Systems that fail any daily verification check shall be placed in
  maintenance status and shall not be used for patient procedures until
  the issue is resolved and the check is passed.

\item Daily check results shall be recorded in the system's audit trail with
  timestamp and pass or fail status for each verification element.
\end{enumerate}

\subsection{System Utilization Management}

\begin{enumerate}[label=(\alph*)]
\item The site's scheduling system shall manage Physical AI system utilization
  to balance patient demand against system availability, following the
  utilization patterns validated in the simulation:
  \begin{enumerate}[label=(\arabic*)]
  \item RT positioning robots: 38 percent average utilization.
  \item Imaging assistants: 38 percent average utilization.
  \item RT motion-tracking robots: 37 percent average utilization.
  \item Surgical robots: scheduled with appropriate intervals between
    procedures for decontamination and preparation.
  \item Companion robots and humanoids: flexible scheduling to accommodate
    patient emotional support needs~\cite{newtrial2026}.
  \end{enumerate}

\item When utilization approaches 80 percent for any system type, the
  scheduling system shall extend appointment intervals to prevent queue
  buildup and maintain the target 8-minute average wait time.

\item Maintenance windows shall be scheduled during periods of historically
  low demand (typically 22:00 to 06:00), using the simulation data as the
  baseline demand model~\cite{newtrial2026}.
\end{enumerate}

\section{Quality Assurance}

\subsection{Continuous Quality Monitoring}

\begin{enumerate}[label=(\alph*)]
\item The site shall maintain a quality management system compliant with the
  adapted ICH E6(R3) requirements for Physical AI
  trials~\cite{iche6r3adapt}.

\item Key quality indicators shall be tracked continuously:
  \begin{enumerate}[label=(\arabic*)]
  \item Patient throughput versus the 168-patient 24-hour baseline.
  \item Adverse event rate versus the 4.2 percent baseline (7 events across
    168 patients)~\cite{newtrial2026}.
  \item Same-hour adverse event resolution rate versus the 100 percent
    baseline.
  \item Fleet uptime versus the 99.7 percent baseline.
  \item Patient satisfaction versus the 4.7 out of 5.0 baseline.
  \item Cumulative PSL score versus the 63.4 to 64.4 baseline range.
  \end{enumerate}

\item Deviations from baseline quality indicators shall trigger root cause
  analysis and corrective action within 7 days.
\end{enumerate}

\subsection{Regulatory Compliance Monitoring}

\begin{enumerate}[label=(\alph*)]
\item Compliance with the three adapted regulatory frameworks shall be
  monitored continuously:
  \begin{enumerate}[label=(\arabic*)]
  \item ICH E6(R3): GCP compliance for Physical AI systems, data governance,
    audit trail integrity~\cite{iche6r3adapt}.
  \item 21 CFR Part 50: Informed consent completeness, patient rights
    adherence, pediatric protections~\cite{cfr50adapt}.
  \item 21 CFR Part 312: IND compliance, safety reporting timelines, annual
    report preparation~\cite{cfr312adapt}.
  \end{enumerate}

\item Monthly compliance reports shall be submitted to CDPH as required by
  the state regulations.

\item The Physical AI Safety Review Committee shall meet at intervals not
  exceeding 90 days to review all safety data, quality metrics, and
  compliance status~\cite{cfr312adapt}.
\end{enumerate}

\section{Adverse Event Management}

\begin{enumerate}[label=(\alph*)]
\item Adverse event detection, documentation, and initial response shall be
  automated by Physical AI systems, with site safety officer notification
  within 1 minute of detection.

\item The seven adverse event categories defined in the adapted 21 CFR Part 312
  shall be used for classification: serious Physical AI adverse events,
  Physical AI-drug interactions, cybersecurity incidents, AI model degradation,
  sim-to-real divergence, digital twin discrepancy, and safety committee
  findings~\cite{cfr312adapt}.

\item The site shall target same-hour resolution for all Grade 1 and Grade 2
  events, matching the simulation performance where all seven adverse events
  (nausea Grade 1, hypotension Grade 1, bleeding Grade 1, pain Grade 2,
  cough Grade 1, desaturation Grade 1, anxiety Grade 1) were resolved within
  the hour of occurrence~\cite{newtrial2026}.

\item Grade 3 or higher events shall trigger immediate escalation to the
  on-call emergency physician and the Physical AI Safety Review Committee.

\item All adverse events shall be reported to CDPH and FDA within the timelines
  specified in the adapted 21 CFR Part 312 \S~312.32 (7 days for fatal or
  life-threatening, 15 days for serious)~\cite{cfr312adapt}.
\end{enumerate}


% ============================================================================
% DOCUMENT 11: Physical AI Oncology Clinical Trial Site Emergency Preparedness Plan
% ============================================================================

\part{Physical AI Oncology Clinical Trial Site Emergency Preparedness Plan}
\setcounter{section}{0}

\section{Emergency Classification System}

Physical AI oncology clinical trial sites shall classify emergencies into four
levels, each with defined response protocols:

\subsection{Level 1 -- Minor Incident}

\begin{enumerate}[label=(\alph*)]
\item Definition: An event affecting a single Physical AI system or a single
  patient without immediate risk of serious harm. Examples include a single
  robot entering fault mode, a Grade 1 adverse event, or a minor network
  disruption affecting non-critical systems.

\item Response: The affected Physical AI system is automatically placed in safe
  mode. The site safety officer assesses the situation and determines whether
  the system can be restarted or requires maintenance. Patient care continues
  using alternative systems. The simulation demonstrated that the fleet of 29
  robot instances provides redundancy within each system category, allowing
  continued operations when individual systems are
  unavailable~\cite{newtrial2026}.

\item Reporting: Documented in the shift log and included in the monthly report
  to CDPH.
\end{enumerate}

\subsection{Level 2 -- Moderate Emergency}

\begin{enumerate}[label=(\alph*)]
\item Definition: An event affecting multiple Physical AI systems, multiple
  patients, or requiring activation of emergency medical response. Examples
  include multiple simultaneous robot faults, a Grade 2 or Grade 3 adverse
  event, a localized power failure, or a cybersecurity anomaly.

\item Response: The site safety officer activates the emergency response team.
  The on-call emergency physician is notified. Affected clinical areas are
  secured. Patient procedures in unaffected areas may continue. All active
  procedures in affected areas are safely concluded or aborted using the
  emergency stop protocols described in the patient robot
  instructions~\cite{patientinstructions2026}.

\item Reporting: Reported to CDPH within 72 hours. If a cybersecurity incident
  with potential patient harm, reported within 7 days per the adapted 21 CFR
  Part 312 \S~312.32~\cite{cfr312adapt}.
\end{enumerate}

\subsection{Level 3 -- Major Emergency}

\begin{enumerate}[label=(\alph*)]
\item Definition: An event requiring partial or complete facility evacuation,
  affecting the majority of Physical AI systems, or involving a serious
  adverse event (Grade 4 or higher). Examples include building-wide power
  failure exceeding UPS capacity, structural damage, confirmed cybersecurity
  breach affecting patient safety, or a mass casualty event.

\item Response: The site administrator activates the facility emergency plan.
  All Physical AI systems execute automated safe-shutdown procedures. Patient
  evacuation begins per the evacuation plan. External emergency services are
  activated. The Physical AI Safety Review Committee is convened within 24
  hours~\cite{cfr312adapt}.

\item Reporting: Reported to CDPH and FDA immediately. Fatal or
  life-threatening events reported within 7 days per \S~312.32. Serious
  events reported within 15 days~\cite{cfr312adapt}.
\end{enumerate}

\subsection{Level 4 -- Catastrophic Event}

\begin{enumerate}[label=(\alph*)]
\item Definition: An event rendering the facility completely non-operational,
  such as a major earthquake, widespread infrastructure failure, or complete
  cybersecurity compromise.

\item Response: Complete facility evacuation. All Physical AI systems
  de-energized. Business continuity plan activated. Patient care transferred
  to partner facilities. IND sponsor notified for potential clinical
  hold assessment~\cite{cfr312adapt}.

\item Reporting: Immediate notification to CDPH, FDA, and local emergency
  management. Full incident report within 30 days.
\end{enumerate}

\section{Physical AI System Emergency Procedures}

\subsection{Emergency Stop Protocol}

\begin{enumerate}[label=(\alph*)]
\item Every Physical AI system shall have a clearly marked, accessible emergency
  stop mechanism that can be activated by site safety officers, clinical
  personnel, or patients where clinically appropriate.

\item Upon emergency stop activation, the system shall:
  \begin{enumerate}[label=(\arabic*)]
  \item Immediately cease all autonomous motion.
  \item Retract any instruments or probes that are in contact with or
    inside the patient, using a controlled withdrawal sequence.
  \item Lower treatment couches to the lowest safe position.
  \item De-energize radiation sources (RT systems).
  \item Maintain essential monitoring and life safety functions.
  \item Generate an emergency stop event record in the audit trail with
    timestamp, trigger source, and system state at the time of
    activation~\cite{iche6r3adapt}.
  \end{enumerate}

\item The patient robot instructions for each system type include specific
  emergency behavior; for example, surgical robots retract arms through chest
  ports, RT positioning systems lower couches to exit height, and
  exoskeletons return to standing rest position~\cite{patientinstructions2026}.
\end{enumerate}

\subsection{Cybersecurity Incident Response}

\begin{enumerate}[label=(\alph*)]
\item The cybersecurity incident response plan shall address Physical
  AI-specific threats:
  \begin{enumerate}[label=(\arabic*)]
  \item Robot control system compromise (unauthorized commands to Physical AI
    systems).
  \item AI model poisoning (manipulation of AI model parameters or training
    data).
  \item Digital twin data manipulation (alteration of patient-specific models).
  \item MCP-PAI server infrastructure attack (compromise of authorization,
    FHIR, DICOM, ledger, or provenance servers).
  \item Network segmentation breach (traffic between isolated network
    zones).
  \end{enumerate}

\item Upon detection of a cybersecurity incident:
  \begin{enumerate}[label=(\arabic*)]
  \item All affected Physical AI systems are immediately isolated from the
    network.
  \item Active procedures on affected systems are safely concluded or aborted.
  \item The cybersecurity operations specialist assesses the scope and
    severity.
  \item Patient data integrity is verified using hash-chained audit trail
    validation~\cite{iche6r3adapt}.
  \item The incident is classified and reported per the adapted 21 CFR
    Part 312 timelines~\cite{cfr312adapt}.
  \end{enumerate}

\item The site shall conduct annual penetration testing and quarterly
  vulnerability assessments of all Physical AI system interfaces and
  MCP-PAI infrastructure.
\end{enumerate}

\section{Natural Disaster Preparedness}

\subsection{Seismic Preparedness}

\begin{enumerate}[label=(\alph*)]
\item The facility shall be designed to Seismic Design Category D or higher
  per ASCE 7 and the California Building Code, appropriate for San Francisco's
  seismic zone.

\item Physical AI systems shall be seismically anchored to their mounting
  surfaces with anchorage rated for the systems' operating weight plus 50
  percent margin.

\item Automated seismic sensors shall trigger Physical AI system safe-mode
  procedures when ground acceleration exceeds 0.1g, including retracting
  surgical instruments, lowering couches, and pausing radiation delivery.

\item Post-earthquake inspection protocols shall verify Physical AI system
  calibration before resuming patient procedures.
\end{enumerate}

\subsection{Power Failure}

\begin{enumerate}[label=(\alph*)]
\item The dual utility feed design provides primary power redundancy. If both
  feeds fail, the automatic transfer switch engages emergency generators
  within 10 seconds.

\item UPS systems bridge the transfer gap, providing 30 minutes for surgical
  suites (sufficient to complete or safely abort any procedure) and 15 minutes
  for server rooms (sufficient for orderly AI system shutdown and data
  preservation).

\item If generator fuel is projected to be exhausted before utility power
  restoration, the site shall begin controlled patient discharge and Physical
  AI system shutdown, prioritizing completion of active procedures.

\item The 72-hour generator fuel capacity provides sufficient time for
  coordinated patient transfer to partner facilities in extended outage
  scenarios.
\end{enumerate}

\section{Clinical Emergency Procedures}

\subsection{Medical Emergency During Physical AI Procedure}

\begin{enumerate}[label=(\alph*)]
\item If a patient experiences a medical emergency during a Physical AI-assisted
  procedure:
  \begin{enumerate}[label=(\arabic*)]
  \item The Physical AI system's vital sign monitoring detects the emergency
    and initiates automated response protocols.
  \item Emergency stop is activated automatically or by the site safety officer.
  \item The on-call emergency physician is notified immediately.
  \item The Physical AI system retracts instruments and positions the patient
    for emergency access.
  \item Resuscitation equipment is deployed from the emergency station located
    within each clinical wing.
  \item The event is documented in real time through the Physical AI system's
    audit trail and the patient's electronic health record.
  \end{enumerate}

\item The simulation demonstrated that real-time vital sign monitoring at
  minute-level resolution enables early detection of adverse events; all
  seven adverse events in the simulation were detected, graded, and managed
  through the automated monitoring
  system~\cite{newtrial2026}.
\end{enumerate}

\subsection{Multiple Simultaneous Adverse Events}

\begin{enumerate}[label=(\alph*)]
\item If multiple adverse events occur simultaneously (which did not occur in
  the simulation but is planned for), the site shall:
  \begin{enumerate}[label=(\arabic*)]
  \item Triage events by severity grade, prioritizing the highest-grade event.
  \item Activate the staffing contingency plan to bring additional personnel
    on-site within 30 minutes.
  \item Suspend new patient intake until the simultaneous events are resolved.
  \item Notify the Physical AI Safety Review Committee if the events suggest
    a systemic issue affecting multiple Physical AI systems or patients.
  \end{enumerate}
\end{enumerate}

\section{Business Continuity}

\begin{enumerate}[label=(\alph*)]
\item The site shall maintain mutual aid agreements with at least two hospital
  or clinical facilities within 30 minutes travel time for patient transfer
  during extended emergencies.

\item Critical Physical AI system configuration data, patient records, and audit
  trails shall be backed up to a geographically remote data center with
  recovery point objective (RPO) of 15 minutes and recovery time objective
  (RTO) of 4 hours.

\item The business continuity plan shall include provisions for resuming
  Physical AI operations following a major emergency, including system
  recalibration, Phase 0 re-validation if systems were physically damaged,
  and regulatory notification of operational restart.

\item Annual business continuity exercises shall test the full emergency
  response chain from emergency detection through patient transfer, Physical
  AI system shutdown, data preservation, and operational restart.
\end{enumerate}

\section{Training and Drills}

\begin{enumerate}[label=(\alph*)]
\item All site personnel shall complete emergency preparedness training as part
  of the 160-hour site safety officer certification program.

\item Emergency drills shall be conducted quarterly, rotating through the
  following scenarios:
  \begin{enumerate}[label=(\arabic*)]
  \item Quarter 1: Physical AI system failure during active procedure
    (emergency stop, patient safety, system recovery).
  \item Quarter 2: Cybersecurity incident (network isolation, data integrity
    verification, incident reporting).
  \item Quarter 3: Building evacuation (Physical AI safe-shutdown, patient
    evacuation, external emergency services coordination).
  \item Quarter 4: Extended power failure (generator operation, controlled
    patient discharge, business continuity activation).
  \end{enumerate}

\item Drill performance shall be documented and reviewed by the Physical AI
  Safety Review Committee, with corrective actions for any deficiencies
  identified.

\item The patient journey demonstrated that 10 distinct stages of care, from
  pre-screening through data lock, require coordinated emergency response
  capabilities at each stage~\cite{patientjourney2026}.
\end{enumerate}

\printbibliography[title={References}]

\end{document}
